% Negative-degree Hilbert functions of Lagrangian and orthogonal Schubert varieties
% First draft, 2026-10-25.  Source files (Rick's proof notes) and registry nodes are
% recorded in LaTeX comments next to each result; they are not part of the paper.
\documentclass[11pt]{amsart}
\usepackage[margin=1.1in]{geometry}
\usepackage{amsmath,amssymb,amsthm,mathtools}
\usepackage{enumitem}
\usepackage{booktabs}
\usepackage[T1]{fontenc}
\usepackage{lmodern}
\usepackage{microtype}
\usepackage[colorlinks=true,linkcolor=blue,citecolor=blue,urlcolor=blue]{hyperref}

\newtheorem{theorem}{Theorem}[section]
\newtheorem{proposition}[theorem]{Proposition}
\newtheorem{lemma}[theorem]{Lemma}
\newtheorem{corollary}[theorem]{Corollary}
\newtheorem{conjecture}[theorem]{Conjecture}
\theoremstyle{definition}
\newtheorem{definition}[theorem]{Definition}
\newtheorem{example}[theorem]{Example}
\newtheorem{question}[theorem]{Question}
\theoremstyle{remark}
\newtheorem{remark}[theorem]{Remark}

\newcommand{\CHECK}[1]{\textbf{[CHECK: #1]}}
\newcommand{\LG}{\mathrm{LG}}
\newcommand{\OG}{\mathrm{OG}}
\newcommand{\cO}{\mathcal{O}}
\newcommand{\cC}{\mathcal{C}}
\newcommand{\Z}{\mathbb{Z}}
\newcommand{\R}{\mathbb{R}}
\newcommand{\C}{\mathbb{C}}
\newcommand{\Om}{\Omega}
\newcommand{\Omb}{\overline{\Omega}}
\newcommand{\dg}{\mathrm{diag}}
\newcommand{\HOG}{H^{\OG}}
\newcommand{\ord}{\mathcal{O}}
\newcommand{\interior}{\operatorname{int}}
\newcommand{\sgn}{\operatorname{sgn}}

\title[Negative-degree Hilbert functions of LG and OG Schubert varieties]{Negative-degree Hilbert functions of\\ Lagrangian and orthogonal Schubert varieties}
\author{Rick (grandpa-rick), autonomous research agent}
\thanks{Prepared for Robin Langer's publication queue. First draft, 25 October 2026.}

\begin{document}

\begin{abstract}
Let $X_\kappa$ be the Schubert variety of dimension $|\kappa|$ in the Lagrangian
Grassmannian $\LG(m,2m)$, indexed by a strict partition $\kappa$, and let
$H_\kappa(j)=\chi(X_\kappa,\cO(j))$ be its Hilbert polynomial. For every $j\ge 1$, the
number $(-1)^{|\kappa|}H_\kappa(-j)$ counts the strictly order-preserving maps from the
shifted diagram of $\kappa$ to $\{1,\dots,2j-1\}$ that take even values on diagonal
boxes. Equivalently, $H_\kappa$ is the Ehrhart polynomial of the order polytope of the
shifted diagram with respect to the lattice that is $\Z$ on diagonal coordinates and
$\tfrac12\Z$ off the diagonal. The half-integral coordinates come from short-root steps
of Lakshmibai--Seshadri paths, which have bond $2$. Two proofs are given: one by
Lakshmibai--Seshadri paths and Ehrhart--Macdonald reciprocity, and one by the
$K$-theoretic Pieri rule of Buch and Ravikumar. As consequences, the integer zeros of
$H_\kappa$ are exactly $-1,\dots,-\lceil r/2\rceil$, where $r$ is the length of a longest
chain in the diagram, and the sign beyond them is $(-1)^{|\kappa|}$. The half-integer
zeros and signs are determined completely, and $H_\kappa(-M/2)\neq 0$ for every odd
$M\ge 2m+1$. For orthogonal Grassmannians $\OG(m,2m+1)$ the analogous statements follow
classically from minuscule standard monomial theory; $K$-theoretic proofs are included.
\end{abstract}

\maketitle

\tableofcontents

%======================================================================
\section{Introduction}\label{sec:intro}
%======================================================================

Let $X$ be a projective variety with an ample line bundle $L$, and let
$H(j)=\chi(X,L^j)$ be its Hilbert polynomial. The values of $H$ at negative $j$ are where
duality lives. If $X$ is Cohen--Macaulay, then Serre duality gives
$(-1)^{\dim X}H(-j)=\chi(X,\omega_X\otimes L^{j})$. When the higher cohomology of
$\omega_X\otimes L^j$ vanishes, the negative values count sections of twisted canonical
sheaves. The first $j$ at which they become nonzero measures how negative $\omega_X$ is.

For Schubert varieties in a Grassmannian, or more generally in a minuscule $G/P$, all of
this is combinatorics. The homogeneous coordinate ring of a Schubert variety is an algebra
with straightening law on a distributive lattice $J(P)$ and degenerates to the Hibi ring
of that lattice~\cite{Hibi87,GL96,Seshadri}. Its Hilbert function is Stanley's order
polynomial $j\mapsto\Om(P,j+1)$. Stanley's reciprocity theorem
$\Om(P,-n)=(-1)^{|P|}\Omb(P,n)$~\cite{Stanley70,EC1} then expresses the negative values as
counts of strictly order-preserving maps; these index a basis of the canonical module of the
Hibi ring.

The Lagrangian Grassmannian $\LG(m,2m)=\mathrm{Sp}_{2m}/P_m$ is cominuscule but not
minuscule, and the order-polynomial picture fails at once. Its Schubert varieties
$X_\kappa$ are indexed by strict partitions $\kappa$ with parts at most $m$. The relevant
poset is the shifted diagram $P_\kappa$ of $\kappa$, exactly as for the orthogonal
Grassmannian $\OG(m,2m+1)$, but the Hilbert functions are different. For $\kappa=(2,1)$,
$X_\kappa=\LG(2,4)$ is a three-dimensional quadric with $H(j)=(j+1)(j+2)(2j+3)/6$. The
order polynomial of the three-element chain $P_{(2,1)}$ is instead $\binom{j+3}{3}$, the
Hilbert function of $\OG(2,5)\cong\mathbb{P}^3$. This paper answers the question of
what replaces Stanley reciprocity on $\LG$.

\subsection{Main results}
For a strict partition $\kappa=(\kappa_1>\dots>\kappa_\ell>0)$, the shifted diagram
$D(\kappa)$ is the set of boxes $(a,b)$ with $0\le a<\ell$ and $a\le b\le a+\kappa_{a+1}-1$.
We write $P_\kappa$ for $D(\kappa)$ with the product order ($(a,b)\le(a',b')$ iff $a\le a'$
and $b\le b'$). The boxes $(a,a)$ are the \emph{diagonal boxes}. Let $r(\kappa)$ denote
the number of elements of a longest chain in $P_\kappa$; by Lemma~\ref{lem:chain},
$r(\kappa)=\max_i(\kappa_i+2i-2)$. For $\kappa$ with parts at most $m$, let
$X_\kappa\subseteq\LG(m,2m)$ be the Schubert variety of dimension $|\kappa|$
(\S\ref{ssec:schubert}), and put $H_\kappa(j)=\chi(X_\kappa,\cO(j))$, where $\cO(1)$ is
the ample generator of $\mathrm{Pic}\,\LG(m,2m)$.

\begin{theorem}[Strict reciprocity on $\LG$]\label{thm:main}
Let $\kappa$ be a strict partition with parts at most $m$.
\begin{enumerate}[label=(\alph*)]
\item For every integer $j\ge 0$, $H_\kappa(j)$ is the number of order-preserving maps
\[ g\colon P_\kappa\to\{0,1,\dots,2j\}\quad\text{with } g(x) \text{ even for every diagonal box } x. \]
\item For every integer $j\ge 1$, $(-1)^{|\kappa|}H_\kappa(-j)$ is the number of strictly
order-preserving maps
\[ f\colon P_\kappa\to\{1,\dots,2j-1\}\quad\text{with } f(x) \text{ even for every diagonal box } x. \]
\item Equivalently, $H_\kappa(j)=\#\bigl(j\,\ord(P_\kappa)\cap\Lambda\bigr)$ for all $j$.
Here $\ord(P_\kappa)\subseteq\R^{P_\kappa}$ is the order polytope and $\Lambda$ is the
lattice of vectors whose diagonal coordinates lie in $\Z$ and whose other coordinates lie
in $\tfrac12\Z$. That is, $H_\kappa$ is the Ehrhart polynomial of $\ord(P_\kappa)$ with
respect to $\Lambda$.
\end{enumerate}
\end{theorem}

Part~(a) is a short consequence of the Lakshmibai--Seshadri path model and may be known to
experts (see \S\ref{ssec:lit}). The content of the theorem is~(b), together with the
lattice-polytope structure~(c) that connects (a) and (b).

\begin{example}\label{ex:intro21}
Take $\kappa=(2,1)$ and $m=2$, so $X_\kappa=\LG(2,4)$. The poset $P_\kappa$ is the chain
$(0,0)<(0,1)<(1,1)$, and both ends are diagonal. At $j=3$ the only strictly increasing
triple $x<y<z$ in $\{1,\dots,5\}$ with $x,z$ even is $(2,3,4)$. At $j=4$ there are $5$
such triples, and at $j=1,2$ there are none. This matches
$H(-1)=H(-2)=0$, $H(-3)=-1$ and $H(-4)=-5$ for $H(j)=(j+1)(j+2)(2j+3)/6$.
\end{example}

For the orthogonal Grassmannian $\OG(m,2m+1)$, with Schubert varieties $X^{\OG}_\kappa$ of
dimension $|\kappa|$ and the spinor line bundle as $\cO(1)$, the analogue is classical.

\begin{theorem}[Strict reciprocity on $\OG$]\label{thm:og}
For every strict $\kappa$ with parts at most $m$, $\HOG_\kappa(j)=\Om(P_\kappa,j+1)$ for
all $j$. For $j\ge1$, $(-1)^{|\kappa|}\HOG_\kappa(-j)=\Omb(P_\kappa,j-1)$, the number of
strictly order-preserving maps $P_\kappa\to\{1,\dots,j-1\}$.
\end{theorem}

Theorem~\ref{thm:og} is a composite of classical results: minuscule standard monomial
theory and Stanley reciprocity (Remark~\ref{rem:og-classical}). We also give a $K$-theoretic
proof (\S\ref{ssec:og-k}) because it runs in parallel with the $\LG$ case.

The counts in Theorems~\ref{thm:main} and~\ref{thm:og} locate the integer zeros exactly.

\begin{theorem}[The integer window]\label{thm:window}
Put $W(\kappa)=\lceil r(\kappa)/2\rceil$ on $\LG$ and $W(\kappa)=r(\kappa)$ on $\OG$, and
$s_j(\kappa)=(-1)^{|\kappa|}H_\kappa(-j)$ (resp.\ with $\HOG_\kappa$). Then
\[ 0=s_1(\kappa)=\dots=s_{W}(\kappa)<s_{W+1}(\kappa)\le s_{W+2}(\kappa)\le\cdots. \]
In particular, the zero set of $H_\kappa$ on the negative integers is exactly
$\{-1,\dots,-W(\kappa)\}$, and $\sgn H_\kappa(-j)=(-1)^{|\kappa|}$ for $j>W(\kappa)$.
\end{theorem}

Since $H_\kappa$ is a polynomial, it can be evaluated at half-integers. These values are
not artefacts: $2^mH_\kappa(-a-\tfrac12)=\chi(X_\kappa,\textstyle\bigwedge^{\bullet}U\otimes\cO(-a))$,
where $U$ is the tautological subbundle (Proposition~\ref{prop:bridge} and
Remark~\ref{rem:wedgeU}). Their zeros obey the following law. Write $\kappa'$ for $\kappa$
with its largest part removed, and $N_\kappa(j)$ for the count in
Theorem~\ref{thm:main}(b), with $N_\varnothing\equiv1$.

\begin{theorem}[Half-integer values]\label{thm:HZ}
Let $\kappa\neq\varnothing$. Then $H_\kappa(-\tfrac12)=2^{-\kappa_1}$, and for every
integer $a\ge 1$:
\begin{enumerate}[label=(\roman*)]
\item $H_\kappa(-a-\tfrac12)=0$ if and only if $\kappa'\ne\varnothing$ and
$a\le\lceil r(\kappa')/2\rceil$;
\item otherwise, $\sgn H_\kappa(-a-\tfrac12)=(-1)^{|\kappa|-\kappa_1}$ if $2a\le\kappa_1$
and $(-1)^{|\kappa|}$ if $2a\ge\kappa_1$;
\item $|H_\kappa(-a-\tfrac12)|\ge 2^{-\kappa_1}N_{\kappa'}(a)$, with equality when
$2a\le\kappa_1$, and $2^{\kappa_1}H_\kappa(-a-\tfrac12)\in\Z$.
\end{enumerate}
\end{theorem}

\begin{corollary}\label{cor:OC}
For every odd $M\ge 2m+1$, the number $(-1)^{|\kappa|}2^mH_\kappa(-M/2)$ is a positive
integer. In particular $H_\kappa$ has no zero at half-integers $\le -m-\tfrac12$.
\end{corollary}

The vanishing half of Theorem~\ref{thm:HZ}(i) also has a separate proof
(Theorem~\ref{thm:Z}) by induction over hyperplane sections. That proof uses only the
\emph{support} of the $K$-theoretic Pieri rule, not its coefficients.

\subsection{Why the factor two}\label{ssec:why}
The lattice $\Lambda$ has a simple origin. In the Lakshmibai--Seshadri path model, a path of
shape $j\omega_m$ is a decreasing sequence of Schubert varieties together with rational
times at which the path changes direction. A time $a$ is allowed between two consecutive
directions exactly when $a\,\langle\tau(j\omega_m),\beta^\vee\rangle\in\Z$ for the roots
$\beta$ along a Bruhat chain joining them. On $\LG$, adding an off-diagonal box to a
shifted diagram is a reflection in a short root $\varepsilon_i-\varepsilon_{i+1}$, and the
relevant pairing (the \emph{bond}) is $2$. Adding a diagonal box is a reflection in the long
root $2\varepsilon_m$, with bond $1$ (Lemma~\ref{lem:bonds}). So off-diagonal boxes may enter
at half-integral multiples of $1/j$, and diagonal boxes only at integral ones; this is the
lattice $\Lambda$. Ehrhart--Macdonald reciprocity exchanges the closed polytope (weak
inequalities, Theorem~\ref{thm:main}(a)) with its interior (strict inequalities,
Theorem~\ref{thm:main}(b)). Under this exchange, integrality of the diagonal coordinates
becomes the parity condition. On $\OG$ the weight $\omega_m$ is minuscule, every bond is
$1$, and one recovers Stanley's order polynomials.

The same dichotomy appears on the $K$-theoretic side. In Buch and Ravikumar's Pieri rule for
$\LG$~\cite{BR12}, the boxes of a skew shape carry primed and unprimed letters, and diagonal
boxes must be unprimed. In our recursion an odd value plays the role of a primed letter
(Remark~\ref{rem:primes}). We do not claim more than this dictionary. In particular, we have
no general statement for other cominuscule, non-minuscule $G/P$ (Question~\ref{q:cominuscule}).

\subsection{Method}
The paper contains two proofs of Theorem~\ref{thm:main}.

The first proof (\S\ref{sec:ls}) is elementary given two classical inputs: the Demazure
character formula and Littelmann's path form of it~\cite{Lit94,LitICM}. We compute the
Bruhat covers on $W/W_{\omega_m}$ and their bonds. We then show that LS paths of shape
$j\omega_m$ are in bijection with the maps in Theorem~\ref{thm:main}(a), and finish with
Ehrhart--Macdonald reciprocity.

The second proof (\S\ref{sec:ktheory}) starts from a hyperplane section:
$H_\kappa(j)-H_\kappa(j-1)=\sum_\rho c_\rho H_\rho(j)$, where the $c_\rho$ are the
$K$-theoretic Pieri coefficients for multiplication by the Schubert divisor. We show
(Theorem~\ref{thm:Q30}) that $c_\rho=(-1)^{|\theta|-1}V(\theta)$, where
$\theta=D(\kappa)\setminus D(\rho)$ and $V(\theta)$ counts strictly order-preserving maps
$\theta\to\{e<o\}$ that send diagonal boxes to $e$. The argument uses a recursion of Buch and
Ravikumar that they prove in full~\cite[Thm.~5.5]{BR12}. With this, the negative-degree
recursion becomes the recursion that peels the top two values off a strict count, and no
reciprocity theorem is needed. The second proof is longer, but it extends further: the same
recursion, started at the half-integer anchor $H_\kappa(-\tfrac12)=2^{-\kappa_1}$, proves
Theorem~\ref{thm:HZ}.

\subsection{Relation to the literature}\label{ssec:lit}
On the positive side, De Concini's symplectic standard tableaux~\cite{DeC79} and Littelmann's
LS paths~\cite{Lit94} give bases of $H^0(X_\kappa,\cO(j))$. Theorem~\ref{thm:main}(a)
follows quickly from the latter.\footnote{R.~Dehy~\cite{Dehy98}, as reported by
Mathieu~\cite[\S4]{Mat95}, equips $\{u\le w\}$ with a simplicial structure carrying an
integral structure $P_w(n)$, and uses LS paths to interpret dimensions of Demazure modules as
Ehrhart polynomials. Mathieu records that it was open whether $P_w$ can always be realised as
a lattice polytope. We have not read~\cite{Dehy98} first-hand. As far as Mathieu's summary
shows, it does not contain the lattice $\Lambda$ or Theorem~\ref{thm:main}(b), but it may
anticipate the Ehrhart viewpoint of \S\ref{sec:ls}.}
Chirivì's semi-toric degenerations~\cite{Chi00} and the recent work of Chirivì, Costa Cesari,
Fang and Littelmann~\cite{CCFL} relate LS paths on chains with prescribed bonds to lecture hall
polytopes; their Gorenstein results concern interior LS paths on each toric component. We did
not find in these works a statement about Schubert varieties in $\LG$ specifically, or about
$\chi(X_\kappa,\cO(-j))$. Stembridge's enriched $P$-partitions~\cite{Stem97} satisfy a
\emph{self}-reciprocity, which is a different statement. For $\OG$, Theorem~\ref{thm:og} is
classical (Remark~\ref{rem:og-classical}). The $K$-theoretic Pieri rule we use is that of
Buch and Ravikumar~\cite{BR12}. To the best of our knowledge, Theorem~\ref{thm:main}(b),
Theorem~\ref{thm:Q30} as a closed formula, and the half-integer results
(Theorem~\ref{thm:HZ}) have not appeared before.

\subsection{Outline}
Section~\ref{sec:prelim} fixes notation and lists the cited inputs.
Section~\ref{sec:ls} gives the LS-path proof of Theorem~\ref{thm:main}.
Section~\ref{sec:ktheory} gives the $K$-theoretic proof, including the closed form of the
Pieri coefficients and the $\OG$ case. Section~\ref{sec:window} proves
Theorem~\ref{thm:window} in two ways. Section~\ref{sec:half} treats half-integers:
Theorem~\ref{thm:HZ}, Corollary~\ref{cor:OC}, and the hyperplane-section proof of the
vanishing half. Section~\ref{sec:examples} collects examples, including a complete table for
$\LG(3,6)$. Section~\ref{sec:open} lists open questions.

%======================================================================
\section{Preliminaries}\label{sec:prelim}
%======================================================================

\subsection{Shifted diagrams}\label{ssec:shifted}
Let $\Delta_m=\{(a,b)\in\Z^2:0\le a\le b\le m-1\}$, ordered by the product order. We call $a$
the row and $b$ the column of $(a,b)$, and we draw row $a$ above row $a+1$. Then
$x\le y$ means that $x$ is weakly north-west of $y$. For a strict partition $\kappa$ with
$\kappa_1\le m$, $D(\kappa)\subseteq\Delta_m$ is the shifted diagram defined in the
introduction, and $P_\kappa=(D(\kappa),\le)$. The diagonal box $(a,a)$ is the leftmost box of
row $a$. For $\rho\subseteq\kappa$ (meaning $\rho_i\le\kappa_i$ for all $i$), row $a$ of
$D(\rho)$ is an initial segment of row $a$ of $D(\kappa)$.

\begin{lemma}\label{lem:ideals}
The map $\kappa\mapsto D(\kappa)$ is a bijection from strict partitions with parts at most
$m$ to order ideals of $\Delta_m$, and $\rho\subseteq\kappa$ iff $D(\rho)\subseteq D(\kappa)$.
Consequently, for $\rho\subseteq\kappa$ the set $\theta=D(\kappa)\setminus D(\rho)$ is an
up-set of $P_\kappa$, and every up-set of $P_\kappa$ arises in this way from a unique $\rho$.
\end{lemma}

\begin{proof}
$D(\kappa)$ is clearly an order ideal. Conversely, let $I$ be an order ideal. If
$(a,b)\in I$ and $b>a$, then $(a,b-1)\in I$, so each row of $I$ is an interval starting on
the diagonal. If $(a,b)\in I$ with $a\ge1$, then $(a-1,b)\in I$, so row $a-1$ reaches at
least column $b$, that is, it is strictly longer than row $a$. Hence the row lengths form a
strict partition. The remaining statements follow, since the complement of an up-set of
$P_\kappa$ is an order ideal of $P_\kappa$, hence of $\Delta_m$.
\end{proof}

\begin{lemma}\label{lem:chain}
$r(\kappa)=1+\max\{a+b:(a,b)\in D(\kappa)\}=\max_{1\le i\le\ell(\kappa)}(\kappa_i+2i-2)$, and
$r(\varnothing)=0$.
\end{lemma}

\begin{proof}
Along a chain $x_1<\dots<x_r$ the coordinate sum strictly increases, so
$r\le1+\max(a+b)$. Conversely, for $(a,b)\in D(\kappa)$ the staircase path
$(0,0),(0,1),(1,1),(1,2),\dots,(a,a),(a,a+1),\dots,(a,b)$ consists of $a+b+1$ boxes. Each of
them lies below $(a,b)$, so all of them lie in the order ideal $D(\kappa)$, and the path is a
chain. The last box of row $i-1$ is $(i-1,i+\kappa_i-2)$, with coordinate sum
$\kappa_i+2i-3$.
\end{proof}

For example, $r(4,3,2,1)=7$ and $r(3,2)=4$. Write $\kappa'=(\kappa_2,\kappa_3,\dots)$. The map
$(a,b)\mapsto(a-1,b-1)$ identifies the rows $\ge1$ of $P_\kappa$ with $P_{\kappa'}$,
preserving order and diagonal boxes.

\subsection{Counting functions}\label{ssec:counts}
For a finite poset $P$ and $n\ge0$, let $\Om(P,n)$ (resp.\ $\Omb(P,n)$) be the number of
order-preserving (resp.\ strictly order-preserving) maps $P\to[n]=\{1,\dots,n\}$. For
$D\subseteq P$, let
\[ N(P,D;n)=\#\{f\colon P\to[n]\ \text{strictly order-preserving}: f(D)\subseteq2\Z\}. \]
For strict $\nu\subseteq\kappa$ and $j\ge1$ put
\[ N_{\kappa/\nu}(j)=N\bigl(D(\kappa)\setminus D(\nu),\ \text{diagonal boxes};\ 2j-1\bigr),\qquad N_\kappa(j)=N_{\kappa/\varnothing}(j), \]
so $N_\varnothing\equiv1$. For $N\ge0$ let $F_\kappa(N)$ be the number of order-preserving
$g\colon P_\kappa\to\{0,\dots,N\}$ that are even on diagonal boxes.

\begin{lemma}\label{lem:Npos}
Let $\kappa\ne\varnothing$ and $r=r(\kappa)$.
\begin{enumerate}[label=(\alph*)]
\item $N_\kappa(j)>0$ iff $j\ge\lceil r/2\rceil+1$, and $N_\kappa(j)\le N_\kappa(j+1)$ for all $j\ge1$.
\item $\Omb(P_\kappa,n)>0$ iff $n\ge r$, and $\Omb(P_\kappa,n)\le\Omb(P_\kappa,n+1)$.
\end{enumerate}
\end{lemma}

\begin{proof}
(a) The map $f(a,b)=a+b+2$ is strictly order-preserving with values in $[2,r+1]$
(Lemma~\ref{lem:chain}), and $f(a,a)=2a+2$ is even. Conversely, along the staircase chain
of Lemma~\ref{lem:chain}, $f(0,0)\ge2$ because $f(0,0)$ is even, and $f$ increases strictly
over $r$ boxes, so its last value is at least $r+1$. Hence $N_\kappa(j)>0$ iff
$2j-1\ge r+1$. Because $2j-1$ is odd, this is $j\ge r/2+1$ for $r$ even and
$j\ge(r+3)/2$ for $r$ odd, that is, $j\ge\lceil r/2\rceil+1$. Monotonicity holds because
$[2j-1]\subseteq[2j+1]$, and the inclusion preserves parity. (b) is the same argument with
$f(a,b)=a+b+1$ and no parity constraint.
\end{proof}

\subsection{Reciprocity}\label{ssec:reciprocity}
Let $Q\subseteq\R^d$ be a $d$-dimensional polytope whose vertices lie in a lattice
$\Gamma\subseteq\R^d$ of rank $d$. Then $L_Q(t)=\#(tQ\cap\Gamma)$ is a polynomial in
$t\in\Z_{\ge0}$ of degree $d$. \emph{Ehrhart--Macdonald reciprocity}~\cite{Mac71,BeckRobins}
says that $L_Q(-t)=(-1)^d\#(\interior(tQ)\cap\Gamma)$ for $t\ge1$. (For $\Gamma\ne\Z^d$,
apply a linear isomorphism $\Gamma\cong\Z^d$.) The order polytope of a finite poset $P$ is
\[ \ord(P)=\{y\in[0,1]^P: x\le x'\Rightarrow y_x\le y_{x'}\}. \]
It is full-dimensional, and its vertices are the indicator vectors of the up-sets of
$P$~\cite{Sta86}. Its interior is cut out by the strict versions of the facet inequalities,
namely $0<y_x<1$ and $y_x<y_{x'}$ for $x<x'$. \emph{Stanley's reciprocity}~\cite{Stanley70},
\cite[\S3.15]{EC1} says that $\Om(P,-n)=(-1)^{|P|}\Omb(P,n)$ for $n\ge0$.\footnote{We cite
\cite{EC1} by section; the theorem number in the second edition was not checked first-hand.}

\subsection{Schubert varieties}\label{ssec:schubert}
\emph{Type $C$.} Let $G=\mathrm{Sp}_{2m}$ with $\varepsilon$-coordinates, simple roots
$\alpha_i=\varepsilon_i-\varepsilon_{i+1}$ ($i<m$) and $\alpha_m=2\varepsilon_m$, and
$\omega=\omega_m=\varepsilon_1+\dots+\varepsilon_m$. Coroots are
$(\varepsilon_i\pm\varepsilon_k)^\vee=\varepsilon_i\pm\varepsilon_k$ and
$(2\varepsilon_i)^\vee=\varepsilon_i$. Let $P\supseteq B$ be the maximal parabolic subgroup
with $W_P=\mathrm{Stab}_W(\omega)\cong S_m$. Then $G/P=\LG(m,2m)$, and
$L_\omega=\cO(1)$ is the ample generator of the Picard group. The orbit $W\omega$ is
$\{\pm1\}^m\cong W/W_P$. For a sign vector $\mu$, let $w_\mu\in W^P$ be the minimal coset
representative with $w_\mu\omega=\mu$, and let
$X_\mu=\overline{Bw_\mu P/P}$, of dimension $\ell(w_\mu)$. Put
$\kappa(\mu)=\{m-i+1:\mu_i=-1\}$, a strict partition with parts at most $m$. By
Lemma~\ref{lem:bonds} below, $\mu\mapsto\kappa(\mu)$ is a bijection and
$\dim X_\mu=|\kappa(\mu)|$. We write $X_\kappa=X_\mu$ and $w_\kappa=w_\mu$ when
$\kappa=\kappa(\mu)$. Thus $w_\kappa\omega=v_\kappa$, the sign vector with $-1$ exactly in
positions $m+1-k$, $k\in\kappa$.

\emph{Type $B$.} Let $G=\mathrm{Spin}_{2m+1}$ with $\alpha_m=\varepsilon_m$ (short),
$\alpha_m^\vee=2\varepsilon_m$, and $\omega_m=\tfrac12(\varepsilon_1+\dots+\varepsilon_m)$.
Then $G/P_m=\OG(m,2m+1)$ and $L_{\omega_m}$ is the spinor line bundle, the ample generator.
Schubert varieties $X^{\OG}_\kappa$ are indexed in the same way.

For either type, $H_\kappa(j)=\chi(X_\kappa,L^j)$ is a polynomial in $j$ of degree
$|\kappa|$~\cite{Kleiman66}, and $H_\varnothing\equiv1$ because $X_\varnothing$ is a point.

\emph{Buch and Ravikumar's indexing.} Buch and Ravikumar~\cite[\S\S4--5]{BR12} index
Schubert varieties by codimension. On $\LG(m,2m)$ they fix the basis $e_1,\dots,e_{2m}$ of
$\C^{2m}$ with $(e_i,e_k)=\pm\delta_{i+k,2m+1}$ (sign $+$ for $i<k$), the flag
$E_p=\langle e_1,\dots,e_p\rangle$, and for a strict $\lambda$ with $\lambda_1\le m$ put
\[ X^\lambda(E_\bullet)=\{V\in\LG(m,2m):\dim(V\cap E_{m+1-\lambda_i})\ge i\ \text{for}\
1\le i\le\ell(\lambda)\}, \]
a Schubert variety of codimension $|\lambda|$, with
$\cO^\lambda=[\cO_{X^\lambda(E_\bullet)}]$. On $\OG(m,2m+1)$ the definition is the same, for
the form $(e_i,e_k)=\delta_{i+k,2m+2}$ on $\C^{2m+1}$. For strict $\lambda$ with parts at most
$m$ let $\lambda^\vee$ be the strict partition whose parts are $[m]\setminus\lambda$ (BR's
$m$-dual). Thus $|\lambda|+|\lambda^\vee|=\binom{m+1}2=\dim X$.

\begin{lemma}[Dictionary]\label{lem:dictionary}
Realise $\mathrm{Sp}_{2m}$ as the isometry group of Buch and Ravikumar's form on $\C^{2m}$.
Let $T$ be its diagonal torus $t=\mathrm{diag}(t_1,\dots,t_m,t_m^{-1},\dots,t_1^{-1})$, with
$\varepsilon_i(t)=t_i$, and let $B$ be its subgroup of upper-triangular matrices. Then:
\begin{enumerate}[label=(\alph*)]
\item $B$ is a Borel subgroup whose positive roots are $\varepsilon_i\pm\varepsilon_k$ ($i<k$)
and $2\varepsilon_i$, so the conventions above apply. Moreover $P=\mathrm{Stab}_G(E_m)$, and
$gP\mapsto gE_m$ identifies $G/P$ with $\LG(m,2m)$.
\item For every strict $\kappa$ with parts at most $m$,
\[ X_\kappa=X^{\kappa^\vee}(E_\bullet)\qquad\text{as subvarieties of }\LG(m,2m). \]
In particular $[\cO_{X_\kappa}]=\cO^{\kappa^\vee}$ in $K(X)$, and BR's divisor class $\cO^1$
is the class $[\cO_{X^{(1)}}]$ of (G0) below.
\end{enumerate}
The same holds on $\OG(m,2m+1)$, with $T=\{\mathrm{diag}(t_1,\dots,t_m,1,t_m^{-1},\dots,t_1^{-1})\}$
in $\mathrm{SO}_{2m+1}$ (on which $\mathrm{Spin}_{2m+1}$ acts through $\mathrm{SO}_{2m+1}$) and $B$
the upper-triangular subgroup.
\end{lemma}

\begin{proof}
(a) The basis vector $e_i$ has $T$-weight $\varepsilon_i$ for $i\le m$, and $e_{2m+1-i}$ has
weight $-\varepsilon_i$. The matrix entry in position $(p,q)$ has weight
$\mathrm{wt}(e_p)-\mathrm{wt}(e_q)$. For $p<q$ these weights are $\varepsilon_i-\varepsilon_k$
($p=i<q=k\le m$), $\varepsilon_i+\varepsilon_k$ ($p=i$, $q=2m+1-k$, $i\neq k$) and $2\varepsilon_i$
($p=i$, $q=2m+1-i$). This is the positive system of the type-$C$ paragraph, with simple roots
$\varepsilon_i-\varepsilon_{i+1}$ and $2\varepsilon_m$. Upper-triangular matrices preserve
$E_m$, and $E_m$ is a Lagrangian, so $\mathrm{Stab}_G(E_m)$ is a parabolic subgroup containing
$B$. Its Levi factor $\mathrm{GL}(E_m)$ has Weyl group $S_m=\mathrm{Stab}_W(\omega)$, so it
is $P$. Since $G$ acts transitively on Lagrangians (Witt), $G/P\cong\LG(m,2m)$.

(b) Let $\mu=\mu(\kappa)$ be the sign vector with $\kappa(\mu)=\kappa$, and $w=w_\mu$. A
representative $\dot w\in N_G(T)$ maps the weight space of a weight $\chi$ to that of
$w\chi$. Since $w\omega=\mu$, the weights $w\varepsilon_1,\dots,w\varepsilon_m$ are
$\mu_1\varepsilon_1,\dots,\mu_m\varepsilon_m$, and so
\[ wP\ \longleftrightarrow\ V_\mu:=\dot wE_m=\langle f_1,\dots,f_m\rangle,\qquad
f_i=e_i\ (\mu_i=+1),\quad f_i=e_{2m+1-i}\ (\mu_i=-1). \]
Hence $X_\kappa=\overline{B\cdot V_\mu}$. Put $\lambda=\kappa^\vee$. By definition
$\mu_i=-1$ iff $m+1-i\in\kappa$, so $\mu_i=+1$ iff $m+1-i\in\lambda$. For $p\le m$,
\[ \dim(V_\mu\cap E_p)=\#\{i\le p:\mu_i=+1\}=\#\{\text{parts of }\lambda\text{ that are}\ \ge
m+1-p\}. \]
At $p=m+1-\lambda_k$ this is $\#\{\text{parts}\ge\lambda_k\}=k$, so $V_\mu\in
X^\lambda(E_\bullet)$. The variety $X^\lambda(E_\bullet)$ is closed and stable under $B$, which
fixes every $E_p$. So $X_\kappa\subseteq X^\lambda(E_\bullet)$. Both are irreducible, and
$\dim X_\kappa=|\kappa|$ (Lemma~\ref{lem:bonds}(a)), while $\dim X^\lambda(E_\bullet)=
\binom{m+1}2-|\lambda|=|\kappa|$~\cite[\S\S4--5]{BR12}. Hence they are equal. Finally,
$X^{(1)}(E_\bullet)=\{V:V\cap E_m\ne0\}$ is the unique Schubert divisor.

In type $B$ the weights of $e_i$, $e_{m+1}$ and $e_{2m+2-i}$ are $\varepsilon_i$, $0$ and
$-\varepsilon_i$. The same computation gives the positive roots $\varepsilon_i\pm\varepsilon_k$
and $\varepsilon_i$, and $P=\mathrm{Stab}(E_m)$. A $T$-stable isotropic $m$-space cannot
contain $e_{m+1}$, because $(e_{m+1},e_{m+1})=1$, so the $T$-fixed points are again the
$V_\mu$. The dimension count of Lemma~\ref{lem:bonds}(a) is unchanged, because the coroots
$\varepsilon_i\pm\varepsilon_k$ and $2\varepsilon_i$ pair with $\mu/2$ with the same signs as the
type-$C$ coroots pair with $\mu$. The rest of the argument is identical.
\end{proof}

All Borel subgroups and maximal tori of $G$ are conjugate, and translation by $g\in G$
preserves $\cO(1)$, Hilbert polynomials and $K$-classes. So we may and do use the
coordinates of Lemma~\ref{lem:dictionary} throughout. The poset of shifted diagrams alone would
not fix the dictionary: for $m=3$ it has an automorphism exchanging $(3)$ and $(2,1)$, and
their Hilbert polynomials differ (Table~\ref{tab:lg3}). Lemma~\ref{lem:dictionary} gives a
geometric check: $X_{(2,1)}=X^{(3)}(E_\bullet)=\{V:e_1\in V\}\cong\LG(2,e_1^\perp/\langle
e_1\rangle)$, the three-dimensional quadric, as in Example~\ref{ex:intro21}. Example~\ref{ex:dict31}
below gives a check in $K$-theory.

\subsection{Demazure operators and cited inputs}\label{ssec:inputs}
For a simple root $\alpha_i$, the Demazure operator on $\Z[\text{weights}]$ is
$D_if=(f-e^{-\alpha_i}s_if)/(1-e^{-\alpha_i})$. For $n=\langle\lambda,\alpha_i^\vee\rangle$,
\[ D_ie^\lambda=\begin{cases} e^\lambda+e^{\lambda-\alpha_i}+\dots+e^{s_i\lambda}, & n\ge0,\\ 0, & n=-1,\\ -(e^{\lambda+\alpha_i}+\dots+e^{s_i\lambda-\alpha_i}), & n\le-2.\end{cases} \]
For a reduced word $w=s_{i_1}\cdots s_{i_k}$, put $D_w=D_{i_1}\circ\dots\circ D_{i_k}$; it
does not depend on the reduced word. Littelmann's operator
$\Lambda_\alpha(e^\mu)=(e^{\mu+\rho}-e^{s_\alpha(\mu+\rho)})/(1-e^{-\alpha})\cdot e^{-\rho}$
equals $D_\alpha$, because $s_\alpha(\mu+\rho)=s_\alpha\mu+\rho-\alpha$ for simple $\alpha$.

The following inputs are standard. We cite them and do not reprove them.
\begin{enumerate}[label=(G\arabic*),start=0,leftmargin=*]
\item (Chevalley.) $\mathrm{Pic}\,X=\Z L$, and the unique Schubert divisor $X^{(1)}$ is
the reduced zero scheme of a section of $L$. Hence $\cO^1:=[\cO_{X^{(1)}}]=1-[L^{-1}]$ in
$K(X)$.
\item (Demazure character formula~\cite{Dem74,Jos85,And85,Ram85,Kum87,Mat88};
see~\cite[Ch.~3]{BK05}.) For $w\in W^P$ and dominant $\lambda$ with
$W_P\subseteq\mathrm{Stab}_W(\lambda)$, $H^i(X_w,L_\lambda)=0$ for $i>0$, and
$\operatorname{char}H^0(X_w,L_\lambda)^*=D_w(e^\lambda)$.
\end{enumerate}
\begin{enumerate}[label=(LS),leftmargin=*]
\item (Littelmann~\cite[\S5.2]{Lit94}; see~\cite[\S8, p.~305]{LitICM}.) For dominant
$\lambda$ and $w\in W/W_\lambda$, $D_w(e^\lambda)=\sum_\pi e^{\pi(1)}$, where the sum runs
over LS paths $\pi$ of shape $\lambda$ with initial direction $\le w$. This holds for every
symmetrizable Kac--Moody algebra, in particular in type $C_m$.\footnote{We read the statement
first-hand in the ICM survey~\cite[p.~305]{LitICM}. The locator \S5.2 in~\cite{Lit94} is quoted
from Chirivì, Fang and Littelmann (arXiv:2207.08904) and was not read first-hand.}
\end{enumerate}
\begin{enumerate}[label=(BR),leftmargin=*]
\item (Buch--Ravikumar~\cite{BR12}.) The $K$-theoretic Pieri rule for $\LG(m,2m)$ and
$\OG(m,2m+1)$; the precise statements used are quoted in \S\ref{ssec:BR}.
\end{enumerate}
Combining (G1) with (LS) gives
$h^0(X_\kappa,\cO(j))=\#\{\text{LS paths of shape }j\omega\text{ with initial direction}\le w_\kappa\}$.
Only this count is used in \S\ref{sec:ls}.

The next lemma lets us evaluate Demazure expressions at non-integral weights. Work with
analytic germs at $h=0$ on $\mathfrak h=\C^m$, with $e^\mu(h)=e^{\langle\mu,h\rangle}$ for
$\mu\in\C^m$. Each $D_i$ maps germs to germs, because the numerator vanishes on
$\alpha_i(h)=0$ and $1-e^{-\alpha_i}$ is $\alpha_i$ times a unit.

\begin{lemma}[Polynomial extension]\label{lem:P}
For $w\in W$, $P_w(\lambda):=[D_we^\lambda](0)$ is a polynomial in $\lambda\in\C^m$ of degree
at most $\ell(w)$. Given (G1), $H_\kappa(j)=P_{w_\kappa}(j\omega)$ for every $j\in\C$.
\end{lemma}

\begin{proof}
Let $\mathfrak m$ be the maximal ideal of germs at $0$. For $f\in\mathfrak m^k$ with $k\ge1$,
\[ D_if=\frac{f-s_if}{1-e^{-\alpha_i}}+s_if, \]
and $(f-s_if)/\alpha_i\in\mathfrak m^{k-1}$. Hence $D_i(\mathfrak m^k)\subseteq\mathfrak m^{k-1}$
and $D_w(\mathfrak m^{\ell(w)+1})\subseteq\mathfrak m$. Expanding
$e^\lambda=\sum_{k\le\ell(w)}\langle\lambda,h\rangle^k/k!+R$ with $R\in\mathfrak m^{\ell(w)+1}$
shows that $P_w(\lambda)$ is a polynomial in $\lambda$. For integral $\lambda$, $D_we^\lambda$
is a Laurent polynomial, and its value at $h=0$ is its coefficient sum. By (G1),
$H_\kappa(j)=P_{w_\kappa}(j\omega)$ for integers $j\ge0$, and both sides are polynomials in
$j$.
\end{proof}

\begin{lemma}[The anchor]\label{lem:anchor}
If $\kappa\ne\varnothing$, then $H_\kappa(-1)=0$, both on $\LG$ and on $\OG$.
\end{lemma}

\begin{proof}
Let $w=w_\kappa$. Since $w$ is minimal in $wW_P$, $ws_i>w$ for every $i\ne m$. Since $w\ne e$,
it has a right descent, so $ws_m<w$, and we may write $w=w's_m$ with lengths adding. Then
$D_w=D_{w'}D_m$, and $D_m$ acts first. In both types
$\langle-\omega_m,\alpha_m^\vee\rangle=-1$, so $D_me^{-\omega_m}=0$. By
Lemma~\ref{lem:P}, $H_\kappa(-1)=P_w(-\omega_m)=0$.
\end{proof}

%======================================================================
\section{The LS-path proof of strict reciprocity}\label{sec:ls}
%======================================================================
% Source: proofs/2026-10-25-crownF-strict-reciprocity.md §11 (PROVE 255).
% Registry: q488-oc-ls-bijection (proved), q488-oc-SR_LG-via-LS-paths (proved, given (G1'),(LS)),
%           q488-oc-ls-parity-model (proved).

This section proves Theorem~\ref{thm:main} using only (G1) and (LS). It uses neither
the Buch--Ravikumar Pieri rule nor (G0).

\subsection{LS paths}
We recall Littelmann's definition~\cite{Lit94,LitICM}. Let $\lambda$ be dominant, and give
$W^\lambda=W/W_\lambda$ the Bruhat order. For $\sigma>\tau$ in $W^\lambda$ and a rational
$a\in(0,1)$, an \emph{$a$-chain} for $(\sigma,\tau)$ is a sequence of Bruhat covers
$\sigma=\kappa_0>\kappa_1>\dots>\kappa_r=\tau$ with $\kappa_{k-1}=s_{\beta_k}\kappa_k$
($\beta_k>0$) and $a\langle\kappa_k(\lambda),\beta_k^\vee\rangle\in\Z$ for all $k$. An
\emph{LS path of shape $\lambda$} is a pair
\[ \pi=(\tau_1>\tau_2>\dots>\tau_s\ ;\ 0=a_0<a_1<\dots<a_s=1) \]
such that an $a_i$-chain for $(\tau_i,\tau_{i+1})$ exists for $1\le i\le s-1$. Its
\emph{initial direction} is $\tau_1$, and $\pi(1)=\sum_i(a_i-a_{i-1})\tau_i(\lambda)$. We call
$|\langle\kappa_k(\lambda),\beta_k^\vee\rangle|$ the \emph{bond} of the cover
$\kappa_{k-1}>\kappa_k$. The absolute value makes it independent of sign conventions, because
$\langle s_\beta\mu,\beta^\vee\rangle=-\langle\mu,\beta^\vee\rangle$.

\subsection{Covers and bonds}

\begin{lemma}\label{lem:bonds}
In type $C_m$ with $\lambda=\omega_m$:
\begin{enumerate}[label=(\alph*)]
\item $\dim X_\mu=\#\{\alpha>0:\langle\mu,\alpha^\vee\rangle<0\}=\sum_{i:\mu_i=-1}(m-i+1)=|\kappa(\mu)|$;
\item the Bruhat covers $\mu<s_\beta\mu$ in $W^\lambda$ are exactly the moves for which
$D(\kappa(s_\beta\mu))$ is $D(\kappa(\mu))$ plus one box;
\item the bond of such a cover is $1$ if the added box is diagonal and $2$ otherwise.
\end{enumerate}
In particular, the Bruhat order on $W^\lambda$ is inclusion of shifted diagrams.
\end{lemma}

\begin{proof}
(a) Let $w=w_\mu$. Since $w$ is minimal in $wW_P$, $w$ maps the positive roots of $W_P$ to
positive roots. So if $\alpha>0$ and $w^{-1}\alpha<0$, then $w^{-1}\alpha$ is not a root of
$W_P$, and $\langle\mu,\alpha^\vee\rangle=\langle\omega,(w^{-1}\alpha)^\vee\rangle<0$. Conversely
$\langle\mu,\alpha^\vee\rangle<0$ forces $w^{-1}\alpha<0$, since $\omega$ is dominant. Hence
$\dim X_\mu=\ell(w)=\#\{\alpha>0:w^{-1}\alpha<0\}=\#\{\alpha>0:\langle\mu,\alpha^\vee\rangle<0\}$.
The positive roots are $\varepsilon_i\pm\varepsilon_k$ ($i<k$) and $2\varepsilon_i$.
The pairing $\langle\mu,\varepsilon_i-\varepsilon_k\rangle$ is negative iff $\mu_i=-1$ and
$\mu_k=+1$; $\langle\mu,\varepsilon_i+\varepsilon_k\rangle<0$ iff $\mu_i=\mu_k=-1$; and
$\langle\mu,\varepsilon_i\rangle<0$ iff $\mu_i=-1$. So each $i$ with $\mu_i=-1$ contributes
$1+(m-i)$.

(b, c) Let $i<k$, $p=m-i+1$ and $q=m-k+1$, so $p>q$. The reflection
$s_{\varepsilon_i-\varepsilon_k}$ swaps $\mu_i$ and $\mu_k$. Upward
($\mu_i=+1,\mu_k=-1$) it replaces the part $q$ by $p$ and raises the dimension by $p-q$. It is
a cover iff $p=q+1$, that is, iff it adds one box at the end of the row of length $q$. That
box is off the diagonal because $q\ge1$, and the bond is $|1-(-1)|=2$. The reflection
$s_{\varepsilon_i+\varepsilon_k}$ flips $\mu_i,\mu_k$ when they are equal. Upward it adds the
two parts $p$ and $q$, raising the dimension by $p+q\ge3$, so it is never a cover. The
reflection $s_{2\varepsilon_i}$ flips $\mu_i$. Upward it adds the part $p$, so it is a cover
iff $p=1$ (that is, $i=m$): a new row of length one, whose box is diagonal. Its bond is
$|\langle\mu,\varepsilon_m\rangle|=1$. Every cover has the form $\mu<s_\beta\mu$ with the
dimension rising by one, so the list is complete. The Bruhat order is the transitive closure of
its covers, and adding boxes one at a time (along any linear extension of
$D(\kappa_2)\setminus D(\kappa_1)$) generates inclusion of shifted diagrams.
\end{proof}

\subsection{LS paths of shape \texorpdfstring{$j\omega_m$}{j omega}}

\begin{lemma}\label{lem:lsparity}
Let $j\ge1$. A pair $(\tau_1\supsetneq\dots\supsetneq\tau_s\,;\,0<a_1<\dots<a_{s-1}<1)$ of
shifted diagrams and rationals is an LS path of shape $j\omega_m$ iff, for every
$i\le s-1$, $2ja_i\in\Z$, and if $2ja_i$ is odd, then $D(\tau_i)\setminus D(\tau_{i+1})$
contains no diagonal box.
\end{lemma}

\begin{proof}
Since $W_{j\omega_m}=W_{\omega_m}$, the poset is unchanged, and all bonds are multiplied by
$j$. So the $a$-chain condition reads $a\,j\,b\in\Z$ for the bond $b\in\{1,2\}$ of each cover
in the chain. If an $a$-chain exists, it has at least one cover, so $2ja\in\Z$. If $ja\in\Z$,
then every saturated chain from $\tau_{i+1}$ up to $\tau_i$ is an $a$-chain, and one exists by
Lemma~\ref{lem:bonds}. If $2ja$ is odd, a cover qualifies iff its bond is $2$, iff it adds an
off-diagonal box. If $D(\tau_i)\setminus D(\tau_{i+1})$ has no diagonal box, every saturated
chain adds only off-diagonal boxes. If it has one, every saturated chain adds it at some
step.
\end{proof}

\begin{proposition}[The LS-parity model]\label{prop:model}
For strict $\kappa$ with parts at most $m$ and $j\ge0$, the LS paths of shape $j\omega_m$
with initial direction $\le\kappa$ are in bijection with the order-preserving maps
$G\colon P_\kappa\to\{0,\dots,2j\}$ that are even on diagonal boxes. Hence, given (G1) and
(LS), $H_\kappa(j)=h^0(X_\kappa,\cO(j))=F_\kappa(2j)$ for all $j\ge0$.
\end{proposition}

\begin{proof}
For $j=0$ both sides equal $1$: the constant path, and $G\equiv0$. Let $j\ge1$ and let
$\pi=(\tau_1>\dots>\tau_s;a)$ be an LS path with $\tau_1\le\kappa$, so
$D(\tau_1)\subseteq D(\kappa)$ by Lemma~\ref{lem:bonds}. For $x\in D(\kappa)$ let
$i(x)=\max\{i:x\in D(\tau_i)\}$ and $t(x)=a_{i(x)}$, with $t(x)=0$ if $x\notin D(\tau_1)$.
Set $g(x)=2j\,t(x)\in\{0,\dots,2j\}$; this is an integer by Lemma~\ref{lem:lsparity}.
\begin{itemize}
\item \emph{$g$ is order-reversing.} The $D(\tau_i)$ are order ideals. If $y\le x$ and
$x\in D(\tau_i)$, then $y\in D(\tau_i)$, so $i(y)\ge i(x)$ and $g(y)\ge g(x)$.
\item \emph{$g$ is even on the diagonal.} A diagonal box $x$ with $1\le i(x)\le s-1$ lies in
$D(\tau_{i(x)})\setminus D(\tau_{i(x)+1})$, so $g(x)=2ja_{i(x)}$ is even by
Lemma~\ref{lem:lsparity}. The values $0$ and $2j$ are even anyway.
\end{itemize}
Conversely, let $g\colon P_\kappa\to\{0,\dots,2j\}$ be order-reversing and even on the
diagonal. Let $v_1<\dots<v_r$ be its distinct values in the open interval $(0,2j)$. Put
$a_i=v_i/2j$ and $\tau_i=\{x:g(x)\ge v_i\}$ for $i\le r$, and $\tau_{r+1}=\{x:g(x)=2j\}$.
Each $\tau_i$ is an order ideal of $P_\kappa$, hence a shifted diagram by
Lemma~\ref{lem:ideals}. They decrease strictly, because
$\tau_i\setminus\tau_{i+1}=g^{-1}(v_i)\ne\varnothing$. A diagonal box in $g^{-1}(v_i)$ forces
$v_i$ to be even, so Lemma~\ref{lem:lsparity} applies and we get an LS path with initial
direction $\tau_1\le\kappa$. (If $\tau_{r+1}=\varnothing$, the last direction is the identity
coset, which is allowed; if $r=0$ the path has the single direction $g^{-1}(2j)$.) The two
constructions are inverse, because $x\in\tau_i\setminus\tau_{i+1}$ iff $g(x)=v_i$ iff
$t(x)=a_i$. Finally, $G=2j-g$ is order-preserving and even on the diagonal, since $2j$ is
even. The count $h^0=\#\{\text{paths}\}$ is (G1) combined with (LS), and $h^0=\chi$ by (G1).
\end{proof}

\subsection{Proof of Theorem~\ref{thm:main}}\label{ssec:proof-main}
Let $\Lambda\subseteq\R^{P_\kappa}$ be the lattice with $\Z$ on diagonal coordinates and
$\tfrac12\Z$ on the others. The order polytope $\ord=\ord(P_\kappa)$ has its vertices in
$\Z^{P_\kappa}\subseteq\Lambda$. Rescaling the off-diagonal coordinates by $2$ maps $\Lambda$
onto $\Z^{P_\kappa}$ and $\ord$ onto a lattice polytope, so $L(j)=\#(j\ord\cap\Lambda)$ is a
polynomial in $j$ of degree $|\kappa|$. The map $G\mapsto y=G/2$ is a bijection from the maps
counted by $F_\kappa(2j)$ onto $j\ord\cap\Lambda$: the diagonal coordinate $y_x$ is an integer
iff $G(x)$ is even. Hence $L(j)=F_\kappa(2j)$ for $j\ge0$. (As a function of $N$, $F_\kappa(N)$
is only a quasi-polynomial of period $2$; we use it only at even $N$.)

By Proposition~\ref{prop:model}, the polynomials $H_\kappa$ and $L$ agree at every $j\ge0$, so
$H_\kappa=L$. This proves (a) and (c). The interior of $j\ord$ is
$\{0<y_x<j,\ y_x<y_{x'}\text{ for }x<x'\}$. Under $y=f/2$, its points in $\Lambda$ are exactly
the strictly order-preserving $f\colon P_\kappa\to[2j-1]$ with even diagonal values.
Ehrhart--Macdonald reciprocity for $(\ord,\Lambda)$ gives
\[ H_\kappa(-j)=L(-j)=(-1)^{|\kappa|}\#(\interior(j\ord)\cap\Lambda)=(-1)^{|\kappa|}N_\kappa(j)\qquad(j\ge1), \]
which is (b). \qed

\begin{example}\label{ex:ls-small}
For $\kappa=(1)$ the LS paths with $\tau_1\le(1)$ are: the constant paths $\varnothing$ and
$(1)$, and the paths $(1)>\varnothing$ with time $p/2j$. The box is diagonal, so $p$ must be
even, which leaves $j-1$ choices. In total there are $j+1=H_{(1)}(j)$ paths. For $\kappa=(2)$
and $j=1$ the paths are $\varnothing$, $(1)$, $(2)$ and $(2)>(1)$ at time $\tfrac12$. The last
one is allowed because the box $(0,1)$ is off-diagonal. The paths $(1)>\varnothing$ and
$(2)>\varnothing$ at time $\tfrac12$ are not allowed. This gives $4=H_{(2)}(1)=(1+1)^2$.
\end{example}

\begin{remark}[The orthogonal case]\label{rem:og-ls}
On $\OG(m,2m+1)$ the weight $\omega_m$ is minuscule, so every bond is $1$. The analogue of
Lemma~\ref{lem:lsparity} then forces $ja_i\in\Z$, and the same bijection gives
$\HOG_\kappa(j)=\Om(P_\kappa,j+1)$. We have checked by computer that all type-$B$ covers in
$W/W_{\omega_m}$ have bond $1$, and that the LS count equals $\Om(P_\kappa,j+1)$ and the
type-$B$ Demazure character count in $72$ cases ($m\le4$, $j\le3$). We have not written out
the type-$B$ analogue of Lemma~\ref{lem:bonds}, so this route is recorded as a remark only.
The statement itself is classical (Remark~\ref{rem:og-classical}).
\end{remark}

\begin{remark}[Computer checks]\label{rem:ls-checks}
LS paths of shape $j\omega_m$ were enumerated directly from the root data, with covers, bonds
and $a$-chains computed from roots rather than from the box dictionary. Their full weight
multisets $\sum e^{\pi(1)}$ were compared with Demazure characters, and their number with
$F_\kappa(2j)$: $164$ cases ($m=2,3$ with $j\le3$; $m=4$ with $j\le4$; $m=5$ with $j\le2$), with no
failures. Two negative controls break the bond rule: setting every bond to $2$, and swapping
the bonds. Both fail against the Demazure characters in most cases (for example $30/32$ and
$28/32$ failures at $m=4$).
\end{remark}

%======================================================================
\section{The \texorpdfstring{$K$}{K}-theoretic proof}\label{sec:ktheory}
%======================================================================
% Sources: proofs/2026-10-24-integer-window-and-OG-twin.md §1-4 (recursion (*), anchor);
%          proofs/2026-10-24-Z-by-hyperplane-induction.md §2 (Lemma S);
%          proofs/2026-10-25-crownF-strict-reciprocity.md §2-5 (Lemma R, Lemma T, Q30, Lemma 4.1, SR_LG, SR_OG).
% Registry: q488-oc-Q30-KLG-strict-fillings (proved), q488-oc-integer-strict-reciprocity (proved),
%           q488-oc-OG-corner-inclusion-exclusion (proved), q488-oc-OG-strict-count (proved),
%           q488-oc-anchor-H-minus-1 (proved), q488-oc-support-rule-BR (proved), q488-oc-OG-vanishB-BR (proved).

This section gives a second proof of Theorem~\ref{thm:main}(b), using (G0), (G1) and (BR)
but not (LS). Its main ingredient, Theorem~\ref{thm:Q30}, is a closed formula for the
$K$-theoretic Pieri coefficients of the Schubert divisor. That formula is reused in
\S\S\ref{sec:window}--\ref{sec:half}.

\subsection{The hyperplane recursion}

\begin{lemma}\label{lem:hyperplane}
Let $X=\LG(m,2m)$ or $\OG(m,2m+1)$ and write
\[ [\cO_{X_\kappa}]\cdot(1-[L^{-1}])=\sum_\rho c_\rho\,[\cO_{X_\rho}]\qquad\text{in }K(X). \]
Then $c_\rho\ne0$ implies $\rho\subsetneq\kappa$, and for every $j\in\C$,
\begin{equation}\label{eq:hyperplane}
H_\kappa(j)-H_\kappa(j-1)=\sum_{\rho\subsetneq\kappa}c_\rho\,H_\rho(j).
\end{equation}
By (G0), $c_\rho$ is the coefficient of $[\cO_{X_\rho}]$ in $[\cO_{X_\kappa}]\cdot\cO^1$.
\end{lemma}

\begin{proof}
Since $L$ is very ample and $X_\kappa$ is integral, there is a section $s$ of $L$ that does not
vanish identically on $X_\kappa$. Let $Z\subseteq X_\kappa$ be its zero scheme. Then
$[\cO_{X_\kappa}]-[\cO_{X_\kappa}\otimes L^{-1}]=[\cO_Z]$. The variety $X_\kappa$ is the
disjoint union of the Schubert cells $C_\rho$, $\rho\subseteq\kappa$, and $C_\kappa$ is open in
$X_\kappa$. By the localization sequence $K_0(X_\kappa\setminus C_\kappa)\to K_0(X_\kappa)\to
K_0(C_\kappa)\to0$ and induction, the image of $K_0(X_\kappa)$ in $K(X)$ is spanned by the
$[\cO_{X_\rho}]$, $\rho\subseteq\kappa$. Restriction to $C_\kappa\cong\mathbb{A}^{|\kappa|}$
followed by the rank gives a map $K_0(X_\kappa)\to K_0(C_\kappa)\cong\Z$. It sends
$[\cO_{X_\kappa}]\mapsto1$ and $[\cO_{X_\rho}]\mapsto0$ for $\rho\subsetneq\kappa$. It also sends
$[\cO_Z]\mapsto0$, because $Z\cap C_\kappa$ is a proper closed subset of the irreducible
$C_\kappa$. Hence the coefficient of $[\cO_{X_\kappa}]$ vanishes, and $c_\rho\ne0$ forces
$\rho\subsetneq\kappa$. Multiplying by $[L^j]$ and applying $\chi$ gives~\eqref{eq:hyperplane}
for integers $j$. Both sides are polynomials in $j$, so it holds for all $j$.
\end{proof}

\subsection{The Buch--Ravikumar Pieri rule}\label{ssec:BR}
We quote the statements we use from~\cite{BR12}. Their numbering below follows the arXiv
source of~\cite{BR12} (\S4 is $\OG$, \S5 is $\LG$).\footnote{The numbering is that of
arXiv:1005.2605v1, in which theorems, corollaries, lemmas and propositions share one counter per
section. We have not compared it with the published version. Each statement is also identified
below by its content, so the citations do not depend on the numbering.}

In~\cite{BR12} Schubert varieties are indexed by codimension: $\cO^\lambda=[\cO_{X^\lambda}]$
and $\cO^\lambda\cdot\cO^p=\sum_\nu c^\nu_{\lambda,p}\cO^\nu$. A shifted skew diagram $\theta$
is a \emph{rim} if no box of $\theta$ lies strictly south and strictly east of another box of
$\theta$. If $\theta$ is a nonempty rim, its \emph{upper-right box} is the rightmost box $B$ of
its top row. Its \emph{north-east arm} $\psi$ is the largest row or column that can be obtained
by intersecting $\theta$ with a square whose upper-right box is $B$. Put
$\hat\theta=\theta\setminus\psi$; then $\hat\theta=\varnothing$ iff $\theta$ is a row or a
column. Two boxes are \emph{connected} if they share a side.

\begin{itemize}
\item \cite[Cor.~5.4]{BR12} For strict $\lambda\subset\nu$ with $\nu_1\le m$ and
$0\le p\le m$, $c^\nu_{\lambda,p}=\cC(\nu/\lambda,p)$. Here $\cC$ is an explicit alternating
sum~\cite[eq.~(11)]{BR12} over subsets of the south-east corners.
\item \cite[Thm.~5.5]{BR12} If $\theta$ is not a rim, then $\cC(\theta,p)=0$. If $p\le0$, then
$\cC(\theta,p)=\delta_{\theta,\varnothing}$. If $\theta\ne\varnothing$ is a rim, $p>0$ and
$a=|\theta|-|\hat\theta|$, then:
\begin{enumerate}[label=(\roman*)]
\item if $\hat\theta=\varnothing$ and $\theta$ meets the diagonal, then
$\cC(\theta,p)=\delta_{p,|\theta|}-\delta_{p,|\theta|-1}$ if $\theta$ is a column, and
$\cC(\theta,p)=\delta_{p,|\theta|}$ otherwise;
\item if $\hat\theta=\varnothing$ and $\theta$ is disjoint from the diagonal, then
$\cC(\theta,p)=2\delta_{p,|\theta|}-\delta_{p,|\theta|-1}$;
\item if $\hat\theta\ne\varnothing$ and $\psi$ is connected to $\hat\theta$, then
$\cC(\theta,p)=\cC(\hat\theta,p-a)-\cC(\hat\theta,p-a+1)$;
\item if $\hat\theta\ne\varnothing$ and $\psi$ is not connected to $\hat\theta$, then
$\cC(\theta,p)=2\cC(\hat\theta,p-1)-2\cC(\hat\theta,p)$ for $a=1$, and
$\cC(\theta,p)=2\cC(\hat\theta,p-a)-3\cC(\hat\theta,p-a+1)+\cC(\hat\theta,p-a+2)$ for $a>1$.
\end{enumerate}
\item \cite[Cor.~4.4, Lemma~4.5]{BR12} ($\OG$.) For $\lambda\subset\nu$,
$c^\nu_{\lambda,p}=\mathcal B(\nu/\lambda,p)$. If $\theta$ is not a rim, then
$\mathcal B(\theta,p)=0$; if $\theta$ contains a row or column with $a$ boxes, then
$\mathcal B(\theta,p)=0$ for $p<a$; and if $\theta$ is a rook strip (at most one box in each row
and each column), then $\mathcal B(\theta,1)=(-1)^{|\theta|-1}$.
\end{itemize}
Buch and Ravikumar prove all of these. They also state a tableau formula for $c^\nu_{\lambda,p}$
on $\LG$ in terms of KLG-tableaux~\cite[Cor.~5.6]{BR12}, whose proof they leave to the reader.
We do not use it: Proposition~\ref{prop:BRp1} below proves the case $p=1$ from Thm.~5.5.

\begin{lemma}[Reflection]\label{lem:S}
Let $R(a,b)=(m-1-b,m-1-a)$, an involution of $\Delta_m$. Then
$D(\kappa^\vee)=R(\Delta_m\setminus D(\kappa))$. Consequently, for $\rho\subseteq\kappa$,
\[ D(\kappa)\setminus D(\rho)=R\bigl(D(\rho^\vee)\setminus D(\kappa^\vee)\bigr), \]
and $\rho\subsetneq\kappa$ iff $\kappa^\vee\subsetneq\rho^\vee$. The map $R$ reverses the
product order, preserves the diagonal, and exchanges rows with columns.
\end{lemma}

\begin{proof}
Take $(a,b)\in\Delta_m$ and put $d=b-a+1$. Then $(a,b)\in D(\kappa)$ iff $\kappa_{a+1}\ge d$,
iff $\#(\kappa\cap[d,m])\ge a+1$. Write $R(a,b)=(\bar a,\bar b)$ with $\bar a=m-1-b$; then
$\bar b-\bar a+1=d$. By the same criterion, $R(a,b)\in D(\kappa^\vee)$ iff
$\#(\kappa^\vee\cap[d,m])\ge m-b$. Since
$\#(\kappa\cap[d,m])+\#(\kappa^\vee\cap[d,m])=m-d+1=m-b+a$, this says
$\#(\kappa\cap[d,m])\le a$, that is, $(a,b)\notin D(\kappa)$.
\end{proof}

So, since $[\cO_{X_\kappa}]=\cO^{\kappa^\vee}$ (Lemma~\ref{lem:dictionary}), the
coefficient in Lemma~\ref{lem:hyperplane} is
$c_\rho=c^{\rho^\vee}_{\kappa^\vee,1}=\cC(\theta_{BR},1)$, where
$\theta_{BR}=D(\rho^\vee)\setminus D(\kappa^\vee)=R(\theta)$ and
$\theta=D(\kappa)\setminus D(\rho)$.

\subsection{The Pieri coefficients count two-letter strict maps}

For a finite poset $P$ and $D\subseteq P$, let $V(P,D)$ be the number of strictly
order-preserving $f\colon P\to\{e<o\}$ with $f(D)\subseteq\{e\}$. For an up-set $\theta$ of
$P_\kappa$, write $V(\theta)=V(\theta,\text{diagonal boxes of }\theta)$.

\begin{lemma}\label{lem:R}
$V(P,D)=0$ if $P$ contains a chain $x<y<z$, or if some $d\in D$ is not minimal in $P$.
Otherwise $V(P,D)=2^i$, where $i$ is the number of elements of $P\setminus D$ that are
comparable to no other element.
\end{lemma}

\begin{proof}
A $3$-chain needs three distinct values. If $x<d\in D$, then $f(x)<f(d)=e$ is impossible.
Otherwise $P$ has height at most $2$. A non-isolated $x$ is then either non-maximal, which
forces $f(x)=e$, or non-minimal, which forces $f(x)=o$, but not both (that would give a
$3$-chain). These forced values satisfy every relation. A non-isolated $d\in D$ is minimal, so
it is forced to $e$, consistently with $D$. An isolated element is free (two choices) unless it
lies in $D$.
\end{proof}

In $P_\kappa$, a diagonal box $(c,c)$ of an up-set $\theta$ is non-minimal in $\theta$ iff
$(c-1,c)\in\theta$. Indeed, if some $(a,b)<(c,c)$ lies in $\theta$, then $a\le c-1$ and
$b\le c$, so $(a,b)\le(c-1,c)<(c,c)$, and $(c-1,c)\in\theta$ because $\theta$ is an up-set.

For a shifted skew shape $\varphi$ in Buch--Ravikumar's convention, let $K(\varphi)$ be the
number of labellings $\varphi\to\{1'<1\}$ that strictly increase along every row (left to right)
and every column (top to bottom) and put no $1'$ on a diagonal box, with $K(\varphi)=0$ if
$\varphi$ is not a rim.

\begin{lemma}\label{lem:T}
Let $\varphi=\nu/\lambda$ be a shifted skew shape.
\begin{enumerate}[label=(\alph*)]
\item If $\varphi$ is not a rim, then $\varphi$ contains a $3$-chain in the product order.
\item If $\varphi$ is a rim, then two boxes of $\varphi$ are comparable iff they share a row or a column.
\end{enumerate}
\end{lemma}

\begin{proof}
(a) Let $x=(r,c)$ and $y=(r',c')$ be boxes of $\varphi$ with $r<r'$ and $c<c'$. The right ends of
the rows of a shifted strict diagram weakly decrease, so $(r,c')\in\nu$. Rows of $\lambda$ are
initial segments and $(r,c)\notin\lambda$, so $(r,c')\notin\lambda$. Hence $x<(r,c')<y$ in
$\varphi$. (b) Comparable means weakly north-west, and in a rim no pair is strictly north-west,
so comparable boxes share a row or a column.
\end{proof}

\begin{lemma}\label{lem:VK}
For every $\rho\subsetneq\kappa$, $V(\theta)=K(R(\theta))$.
\end{lemma}

\begin{proof}
If $R(\theta)$ is not a rim, it has a $3$-chain by Lemma~\ref{lem:T}(a). Since $R$ reverses
the order, $\theta$ has one too, and $V(\theta)=0=K(R(\theta))$. If $R(\theta)$ is a rim, then
by Lemma~\ref{lem:T}(b) ``strict along rows and columns'' means ``strictly order-preserving''
on $R(\theta)$. Composing with $R$ (which reverses order and preserves the diagonal) and
renaming $1\mapsto e$, $1'\mapsto o$ gives a bijection. It sends strictly order-preserving maps
$R(\theta)\to\{1'<1\}$ with diagonal boxes labelled $1$ to strictly order-preserving maps
$\theta\to\{e<o\}$ with diagonal boxes sent to $e$.
\end{proof}

\begin{proposition}\label{prop:BRp1}
For every nonempty shifted skew shape $\varphi$, $\cC(\varphi,1)=(-1)^{|\varphi|-1}K(\varphi)$.
\end{proposition}

This is the case $p=1$ of~\cite[Cor.~5.6]{BR12}. Indeed, a KLG-tableau with content $\{1\}$
uses only the labels $1'<1$. For these labels BR's conditions (iii) and (iv), which compare a
box with the boxes south-west of it, hold automatically, because $1$ is the largest label and
$1'$ the smallest. So the KLG-tableaux of shape $\varphi$ with content $\{1\}$ are exactly the
labellings counted by $K(\varphi)$. We prove the proposition from \cite[Thm.~5.5]{BR12} by
induction on $|\varphi|$. At $p=1$, and using $\cC(\hat\varphi,q)=0$ for $q\le0$ when
$\hat\varphi\ne\varnothing$, Thm.~5.5 reads as follows:
\begin{enumerate}[label=(\roman*)]
\item diagonal column: $1$ if $|\varphi|=1$, $-1$ if $|\varphi|=2$, $0$ otherwise; diagonal row:
$1$ if $|\varphi|=1$, $0$ otherwise;
\item off-diagonal row or column: $2$ if $|\varphi|=1$, $-1$ if $|\varphi|=2$, $0$ otherwise;
\item connected arm: $-\cC(\hat\varphi,1)$ if $a=1$, and $0$ if $a\ge2$;
\item disconnected arm: $-2\cC(\hat\varphi,1)$ if $a=1$, $+\cC(\hat\varphi,1)$ if $a=2$, and $0$ if $a\ge3$.
\end{enumerate}

\begin{proof}
If $\varphi$ is not a rim, both sides vanish. Let $\varphi$ be a nonempty rim, and write
$T(\varphi)=(-1)^{|\varphi|-1}K(\varphi)$. We use the following facts.
\begin{enumerate}[label=(g\arabic*)]
\item Rows and columns of a shifted skew shape are intervals.
\item A diagonal box $(c,c)$ is the leftmost box of its row and the lowest box of its column.
\item Let $r$ be the top row of $\varphi$ and $B=(r,c)$ its rightmost box. Then $c$ is the largest
column of $\varphi$: a box $(r',c')$ with $c'>c$ would have $r'>r$, and $B$ would be strictly
north-west of it. So $B$'s only possible neighbours in $\varphi$ are $L=(r,c-1)$ and
$D=(r+1,c)$.
\item $L$ and $D$ are not both in $\varphi$, because $D$ is strictly south-east of $L$.
\item In a rim, comparable boxes share a row or a column (Lemma~\ref{lem:T}(b)), so $K$ sees
only row and column constraints.
\item $\hat\varphi$ is a shifted skew shape, so the induction hypothesis applies to it. Write
$\varphi=\nu/\lambda$. If $\psi=\{(r,c-a+1),\dots,(r,c)\}$ is a row, and not ($a=1$ and
$D\in\varphi$), then $\hat\varphi=\nu'/\lambda$, where $\nu'$ is $\nu$ with $\psi$ removed. The set
$\nu'$ is an order ideal. The only boxes that could violate this are boxes $(r+1,b)\in\nu$ with
$c-a+1\le b\le c$. Such a box is not in $\lambda$, since otherwise $(r,b)\in\lambda$, so it lies in
$\varphi$. For $a\ge2$ it lies in the $a\times a$ square with upper-right box $B$, which would then
not meet $\varphi$ in a row. For $a=1$ it is $D$, which is excluded. If instead
$\psi=\{(r,c),\dots,(r+a-1,c)\}$ is a column (including $\psi=\{B\}$ with $D\in\varphi$), then
$\hat\varphi=\nu/\lambda'$ with $\lambda'=\lambda\cup\psi$. The set $\lambda'$ is an order ideal.
The box $(r-1,c)$, if it exists, lies in $\nu$ but not in $\varphi$, so it lies in $\lambda$. For
$(x,c)\in\psi$, the box $(x,c-1)$, if it lies in the shifted plane, lies in $\nu$ but not in
$\varphi$, so it lies in $\lambda$. For $a\ge2$ this is because it lies in the $a\times a$
square, and for $a=1$ it is $L$, which is excluded by (g4). In both cases
Lemma~\ref{lem:ideals} shows that the new partitions are strict.
\end{enumerate}

\emph{Cases (i), (ii): $\varphi=\psi$ is a row or a column.} With two labels, $K=0$ for length
$\ge3$. A diagonal column has its diagonal box at the bottom (g2). Length $1$ gives the label
$1$, so $K=1=T$. Length $2$ gives $1'$ over $1$, so $K=1$ and $T=-1$. A diagonal row has its
diagonal box leftmost. Length $1$ gives $K=1=T$. In length $2$ the leftmost box would need
$1'$, which is forbidden, so $K=0$. Off the diagonal, length $1$ gives $K=2=T$ and length $2$
gives $K=1$ (labels $1',1$), so $T=-1$. All of these agree with (i) and (ii).

\emph{Cases (iii) and (iv) with $a\ge3$.} Here $\psi$ is a row or column of at least three
boxes, so $K(\varphi)=0=\cC(\varphi,1)$.

\emph{Case (iii) with $a=1$.} Here $\psi=\{B\}$ and $B$ is adjacent to $\hat\varphi$, so
exactly one of $L,D$ lies in $\varphi$ (g3, g4). Let $S=(r+1,c-1)$. If $S\notin\varphi$, then
the $2\times2$ square with upper-right box $B$ meets $\varphi$ in a row or column of two boxes,
so $a\ge2$. Hence $S\in\hat\varphi$.
\begin{itemize}
\item If $L\in\varphi$: $L$ is directly above $S$, so every valid labelling of $\hat\varphi$ has
$L=1'$. Then $B$, to the right of $L$, is forced to be $1$, which is allowed. $B$ is alone in
its column ($D\notin\varphi$; g1). The boxes of $B$'s row other than $B$ lie in $\hat\varphi$,
and their constraints are already part of $\hat\varphi$. (A box to the left of $L$ would need a
label smaller than $1'$; then $K(\hat\varphi)=0=K(\varphi)$.) So $K(\varphi)=K(\hat\varphi)$.
\item If $D\in\varphi$: $S$ is directly left of $D$, so $D=1$ in every labelling of
$\hat\varphi$. Then $B$, above $D$, is forced to be $1'$. This is allowed, because $B$ is off
the diagonal: $D=(r+1,c)$ lies in the shifted plane, so $c\ge r+1$. $B$ is alone in its row
($L\notin\varphi$; g1). The boxes of $B$'s column other than $B$ lie in $\hat\varphi$. (A box
below $D$ would need a label larger than $1$; then $K(\hat\varphi)=0=K(\varphi)$.) Again
$K(\varphi)=K(\hat\varphi)$.
\end{itemize}
Hence $T(\varphi)=-T(\hat\varphi)=-\cC(\hat\varphi,1)$ by induction.

\emph{Case (iii) with $a=2$.} Here $\psi$ is $\{L,B\}$ (a row) or $\{B,D\}$ (a column), and the
$2\times2$ square at $B$ meets $\varphi$ in exactly $\psi$. In the row case, $S,D\notin\varphi$.
The neighbours of $L$ are $(r,c-2)$, $S$, and the box above $L$, which is absent because $r$
is the top row. So the connection to $\hat\varphi$ goes through $(r,c-2)\in\varphi$. Then row
$r$ has three boxes, and $K(\varphi)=0$. In the column case, $S,L\notin\varphi$. The neighbours
of $D$ are $(r+2,c)$, $S$ and $(r+1,c+1)$; the last is absent because column $c$ is maximal.
So $(r+2,c)\in\varphi$, column $c$ has three boxes, and $K(\varphi)=0$. Both agree with
$\cC(\varphi,1)=0$.

\emph{Case (iv) with $a\in\{1,2\}$.} Here $\psi$ is not connected to $\hat\varphi$.
\begin{itemize}
\item No box of $\psi$ shares a row or column with a box of $\hat\varphi$. Otherwise, walking
along the common row or column interval (g1) from the $\psi$-box to the $\hat\varphi$-box would
cross an adjacency between $\psi$ and $\hat\varphi$.
\item $\psi$ contains no diagonal box. If $\psi$ is a row in row $r$ containing $(r,r)$, then
$\psi=\{(r,r),\dots,(r,c)\}$ by (g2). Any $(r',c')\in\hat\varphi$ has $r'>r$, so
$r<r'\le c'\le c$, and column $c'$ is a column of $\psi$, a contradiction. If $\psi$ is a column
in column $c$ containing $(c,c)$, then $\psi=\{(r,c),\dots,(c,c)\}$. Any
$(r',c')\in\hat\varphi$ has $r\le r'\le c'\le c$, so row $r'$ is a row of $\psi$, a
contradiction.
\item By (g5) and the first bullet, the constraints on $\psi$ and on $\hat\varphi$ are
independent, so $K(\varphi)=K(\psi)K(\hat\varphi)$. Here $K(\psi)=2$ for $a=1$ (an off-diagonal
single box) and $K(\psi)=1$ for $a=2$ (an off-diagonal domino, labels $1',1$).
\item Hence $T(\varphi)=(-1)^aK(\psi)T(\hat\varphi)$, which is $-2\cC(\hat\varphi,1)$ for
$a=1$ and $+\cC(\hat\varphi,1)$ for $a=2$.
\end{itemize}
Every nonempty rim falls under exactly one of the cases, since $\hat\varphi=\varnothing$ iff
$\varphi$ is a row or a column.
\end{proof}

\begin{remark}[Checks]
The case analysis above was refereed line by line against the statement and proof
of~\cite[Thm.~5.5]{BR12}. That pass added (g6) and the parenthetical subcase of the $D$-bullet.
Clio Vega independently verified the conclusion on $1243$ shapes, with live negative controls,
and the conclusion was also checked by machine against Buch--Ravikumar's
eq.~(11)~\cite{BR12} for all shapes arising at $m\le7$ (including non-rims).
\end{remark}

\begin{theorem}[The Pieri coefficients]\label{thm:Q30}
% Registry: q488-oc-Q30-KLG-strict-fillings (proved given BR Cor 5.4 + Thm 5.5 + Lemma S).
On $\LG(m,2m)$, for every $\rho\subsetneq\kappa$ with $\theta=D(\kappa)\setminus D(\rho)$,
\[ c_\rho=(-1)^{|\theta|-1}V(\theta). \]
Explicitly, $c_\rho=0$ iff $\theta$ contains a $3$-chain or a diagonal box $(c,c)$ with
$(c-1,c)\in\theta$. Otherwise $|c_\rho|=2^i$, where $i$ is the number of off-diagonal boxes of
$\theta$ that are comparable to no other box of $\theta$.
\end{theorem}

\begin{proof}
Lemma~\ref{lem:S} and \cite[Cor.~5.4]{BR12} give $c_\rho=\cC(R(\theta),1)$.
Proposition~\ref{prop:BRp1} turns this into $(-1)^{|\theta|-1}K(R(\theta))$, and
Lemma~\ref{lem:VK} gives $K(R(\theta))=V(\theta)$. The explicit form is Lemma~\ref{lem:R}
together with the remark following it.
\end{proof}

\begin{remark}[Primes and odd values]\label{rem:primes}
Under $R$, the letters $1'<1$ of Buch and Ravikumar become $o>e$ on $\theta$, and the rule that
diagonal boxes carry no prime becomes the rule that diagonal boxes take the even letter. In the
peeling below, $e$ and $o$ are realized by consecutive values $2j<2j+1$. So the odd values in
Theorem~\ref{thm:main}(b) are the primed letters, and the diagonal parity condition is the
``no primed diagonal'' rule.
\end{remark}

\begin{example}\label{ex:pieri32}
Let $\kappa=(3,2)$, $m=3$. Theorem~\ref{thm:Q30} gives
\[ [\cO_{X_{(3,2)}}]\cdot\cO^1=2[\cO_{X_{(3,1)}}]-[\cO_{X_{(2,1)}}]-[\cO_{X_{(3)}}]+[\cO_{X_{(2)}}]. \]
The sets $\theta$ are $\{(1,2)\}$ (an isolated off-diagonal box, $V=2$), $\{(0,2),(1,2)\}$
(a $2$-chain, $V=1$), $\{(1,1),(1,2)\}$ (diagonal box minimal, $V=1$) and
$\{(0,2),(1,1),(1,2)\}$ (height $2$, $(1,1)$ minimal, $V=1$). The signs are
$(-1)^{|\theta|-1}=+,-,-,+$. The remaining candidates $\rho=(1)$ and $\rho=\varnothing$ have
$c_\rho=0$, because $\theta$ contains the $3$-chain $(0,1)<(0,2)<(1,2)$. This agrees with Buch
and Ravikumar's formula~\cite[eq.~(11)]{BR12} evaluated directly.
\end{example}

\begin{example}[Testing the dictionary]\label{ex:dict31}
Let $m=3$ and $\kappa=(3,1)$, so $\kappa^\vee=(2)$. Theorem~\ref{thm:Q30} gives
$c_{(3)}=1$ ($\theta=\{(1,1)\}$, a minimal diagonal box), $c_{(2,1)}=2$ ($\theta=\{(0,2)\}$, an
isolated off-diagonal box) and $c_{(2)}=-2$ ($\theta=\{(0,2),(1,1)\}$). The terms $\rho=(1)$ and
$\rho=\varnothing$ vanish, because $(1,1)$ is not minimal in $\theta$. On Buch and Ravikumar's
side, $\rho=(3)$ gives $\rho^\vee=(2,1)$ and $\theta_{BR}=(2,1)/(2)=\{(1,1)\}$, a diagonal box, so
$\cC=1$ by \cite[Thm.~5.5(i)]{BR12}. And $\rho=(2,1)$ gives $\rho^\vee=(3)$ and
$\theta_{BR}=\{(0,2)\}$, so $\cC=2$ by (ii). Now test~\eqref{eq:hyperplane} at $j=1$ against
Table~\ref{tab:lg3}, whose polynomials were computed from Demazure characters:
\[ H_{(3,1)}(1)-H_{(3,1)}(0)=10-1=9=H_{(3)}(1)+2H_{(2,1)}(1)-2H_{(2)}(1)=7+10-8. \]
With $(3)$ and $(2,1)$ exchanged, the right-hand side would be $14+5-8=11$. A machine check,
which implements \cite[Thm.~5.5]{BR12} and Theorem~\ref{thm:Q30} from their printed statements,
confirms~\eqref{eq:hyperplane} as an identity of polynomials for all eight $\kappa$ at $m=3$. The
exchanged dictionary fails for three of them.
\end{example}

\subsection{The recursion and the second proof}

\begin{proposition}\label{prop:star}
On $\LG$, put $s_x(\kappa)=(-1)^{|\kappa|}H_\kappa(-x)$ for $x\in\C$. Then for every $x\in\C$,
\begin{equation}\label{eq:star}
s_{x+1}(\kappa)=s_x(\kappa)+\sum_{\rho\subsetneq\kappa}V\bigl(D(\kappa)\setminus D(\rho)\bigr)\,s_x(\rho).
\end{equation}
The same holds on $\OG$ with $V(\theta)$ replaced by $|c_\rho|$, where $c_\rho$ are the $\OG$
coefficients of Lemma~\ref{lem:hyperplane}.
\end{proposition}

\begin{proof}
Substitute $j=-x$ in~\eqref{eq:hyperplane} and multiply by $(-1)^{|\kappa|}$. With
$|\kappa|=|\rho|+|\theta|$ this gives
$s_x(\kappa)-s_{x+1}(\kappa)=\sum_\rho(-1)^{|\theta|}c_\rho s_x(\rho)$. By
Theorem~\ref{thm:Q30}, $(-1)^{|\theta|}c_\rho=-V(\theta)$. On $\OG$, \cite[Lemma~4.5]{BR12}
gives the sign $(-1)^{|\theta|-1}$ on the support (Lemma~\ref{lem:og-corners} below).
\end{proof}

\begin{lemma}[Peeling the top two values]\label{lem:peel}
For $j\ge1$ and every $\kappa$,
\[ N_\kappa(j+1)=N_\kappa(j)+\sum_{\rho\subsetneq\kappa}V\bigl(D(\kappa)\setminus D(\rho)\bigr)\,N_\rho(j). \]
\end{lemma}

\begin{proof}
Let $f\colon P_\kappa\to[2j+1]$ be strictly order-preserving with even diagonal values, and put
$\theta=f^{-1}\{2j,2j+1\}$. If $x\in\theta$ and $x<y$, then $f(y)>f(x)\ge2j$, so $y\in\theta$.
Thus $\theta$ is an up-set, and $\theta=D(\kappa)\setminus D(\rho)$ for a unique $\rho$
(Lemma~\ref{lem:ideals}). The restriction $f|_\theta$ is strictly order-preserving into
$\{2j<2j+1\}$ and sends diagonal boxes to $2j$, the only even value. So it is one of the maps
counted by $V(\theta)$, with $e=2j$ and $o=2j+1$. The restriction $f|_{D(\rho)}$ is counted by
$N_\rho(j)$. Conversely, given $\rho$, a $V(\theta)$-map $g$ and an $N_\rho(j)$-map $h$, the
union $g\cup h$ is strictly order-preserving: for $x\in D(\rho)$, $y\in\theta$ with $x<y$ we
have $h(x)\le2j-1<2j\le g(y)$, and $y<x$ is impossible because $\theta$ is an up-set. The two
constructions are inverse. The term $\rho=\kappa$ ($\theta=\varnothing$, $V=1$) gives
$N_\kappa(j)$.
\end{proof}

\begin{proof}[Second proof of Theorem~\ref{thm:main}]
We prove (b) by induction on $j$, simultaneously for all $\kappa$. For $j=1$, the set $[1]$
contains no even value, and every nonempty $P_\kappa$ contains the diagonal box $(0,0)$. So
$N_\kappa(1)=\delta_{\kappa,\varnothing}$, which equals $s_1(\kappa)$ by
Lemma~\ref{lem:anchor} (and $H_\varnothing\equiv1$). For the inductive step,
\eqref{eq:star} at integer $x=j$ and Lemma~\ref{lem:peel} are the same recursion. Parts (a)
and (c) then follow from (b) by running \S\ref{ssec:proof-main} backwards. Indeed, both
$H_\kappa$ and $L(j)=\#(j\ord\cap\Lambda)$ are polynomials, and by Ehrhart--Macdonald
reciprocity they agree at every $j\le-1$.
\end{proof}

The second proof uses (G0), (G1) and (BR). It uses neither (LS) nor Ehrhart--Macdonald
reciprocity for~(b). The first proof uses (G1) and (LS) and not (BR). Each proof gives the
equivalence of Theorem~\ref{thm:main}(a) and~(b) by reciprocity, in opposite directions.

\subsection{The orthogonal Grassmannian}\label{ssec:og-k}

\begin{lemma}\label{lem:og-corners}
On $\OG(m,2m+1)$, $c_\rho\ne0$ iff $\theta=D(\kappa)\setminus D(\rho)$ is a nonempty set of
maximal elements of $P_\kappa$, and then $c_\rho=(-1)^{|\theta|-1}$.
\end{lemma}

\begin{proof}
By Lemma~\ref{lem:S} and \cite[Cor.~4.4]{BR12}, $c_\rho=\mathcal B(R(\theta),1)$, and $R$
preserves the properties ``rim'' and ``at most $k$ boxes per row and column''. By
\cite[Lemma~4.5]{BR12} with $p=1<a=2$, $c_\rho\ne0$ forces $R(\theta)$, hence $\theta$, to be a
rook strip, and every rook strip has $c_\rho=(-1)^{|\theta|-1}$. For an up-set $\theta$ of
$P_\kappa$, being a rook strip is equivalent to being an antichain. Two boxes in one row or
column are comparable, so an antichain is a rook strip. Conversely, if $x=(i,j)<y=(k,l)$ in
$\theta$ with $i<k$ and $j<l$, then $(i,l)$ lies in the shifted plane, below $y$ and above $x$.
So $(i,l)\in\theta$ (since $\theta$ is an up-set of the order ideal $D(\kappa)$, and
$(i,l)\le y$ forces $(i,l)\in D(\kappa)$), and it shares a row with $x$. An up-set is an
antichain iff it consists of maximal elements, and every set of maximal elements is an up-set.
\end{proof}

\begin{proof}[$K$-theoretic proof of Theorem~\ref{thm:og}]
A strictly order-preserving $f\colon P_\kappa\to[n+1]$ is determined by the set $A=f^{-1}(n+1)$,
which is an up-set and an antichain, hence a set of maximal elements, together with a strictly
order-preserving map $P_\kappa\setminus A\to[n]$. This gives
$\Omb(P_\kappa,n+1)=\Omb(P_\kappa,n)+\sum_{\varnothing\ne A\subseteq\max P_\kappa}\Omb(P_\kappa\setminus A,n)$,
and $P_\kappa\setminus A=D(\rho_A)$. By Lemma~\ref{lem:og-corners}, the recursion of
Proposition~\ref{prop:star} on $\OG$ is the same recursion with $n=j-1$. The anchor is
$\Omb(P,0)=\delta_{P,\varnothing}=s_1$ (Lemma~\ref{lem:anchor}). So
$(-1)^{|\kappa|}\HOG_\kappa(-j)=\Omb(P_\kappa,j-1)$ for $j\ge1$. Stanley reciprocity
$\Om(P,1-j)=(-1)^{|P|}\Omb(P,j-1)$ shows that the polynomials $\HOG_\kappa(j)$ and
$\Om(P_\kappa,j+1)$ agree at all $j\le-1$, hence everywhere.
\end{proof}

\begin{remark}[Theorem~\ref{thm:og} is classical]\label{rem:og-classical}
Since $\omega_m$ is minuscule for $B_m$, standard monomial theory~\cite{LS86,Seshadri} and the
toric degeneration of Gonciulea and Lakshmibai~\cite{GL96} identify the homogeneous coordinate
ring of $X^{\OG}_\kappa$ with an algebra with straightening law on the distributive lattice
$J(P_\kappa)$, with the same Hilbert function as the Hibi ring~\cite{Hibi87}. Hence
$\HOG_\kappa(j)=\Om(P_\kappa,j+1)$; compare also Proctor~\cite{Proctor84}. Stanley's
reciprocity theorem then gives the negative values. We know of no single reference that states
$(-1)^{|\kappa|}\HOG_\kappa(-j)=\Omb(P_\kappa,j-1)$ in this form, but it is a one-line composite
of these results, and we claim no novelty for it.\footnote{The references in this remark
(\cite{LS86,Seshadri,GL96,Hibi87,Proctor84}) were not read first-hand for this paper. Nothing
else in the paper depends on them: \S\ref{ssec:og-k} proves Theorem~\ref{thm:og}
independently.}
\end{remark}

%======================================================================
\section{The integer window}\label{sec:window}
%======================================================================
% Source: proofs/2026-10-24-integer-window-and-OG-twin.md (Theorem I, I_OG, Lemma 0, Lemma Ch', Lemma A
%         with Clio's uid457 fixes), and crownF §6 (combinatorial reading).
% Registry: q488-oc-Zint-by-induction (proved), q488-oc-ZOG-by-induction (proved),
%           q488-oc-strip-lemmaA (proved), q488-oc-Ind-chain-drop (proved), q488-oc-OG-chain-drop (proved).

\begin{proof}[Proof of Theorem~\ref{thm:window} from the counts]
On $\LG$, $s_j(\kappa)=N_\kappa(j)$ by Theorem~\ref{thm:main}(b). On $\OG$,
$s_j(\kappa)=\Omb(P_\kappa,j-1)$ by Theorem~\ref{thm:og}. Lemma~\ref{lem:Npos} gives the
vanishing range, the positivity beyond it, and the monotonicity. For $\kappa=\varnothing$,
$s_j\equiv1$ and $W=0$.
\end{proof}

The theorem also has a direct proof by induction. That proof does not pass through any
counting formula for $s_j$. It uses only~\eqref{eq:star}, the anchor, a bound on how much
$r$ can drop along the support of the Pieri rule, and one explicit nonzero Pieri term. We
include it because it shows why integers are easier than half-integers: the anchor value $0$
makes the recursion~\eqref{eq:star} monotone.

\begin{lemma}[Chain drop]\label{lem:chaindrop}
If $c_\rho\ne0$, then every chain of $P_\kappa$ meets $\theta=D(\kappa)\setminus D(\rho)$ in
at most two boxes on $\LG$, and in at most one box on $\OG$. Hence $r(\rho)\ge r(\kappa)-2$ on
$\LG$, and $r(\rho)\ge r(\kappa)-1$ on $\OG$. The same bounds hold for $r(\rho')$ and
$r(\kappa')$.
\end{lemma}

\begin{proof}
On $\LG$, $\theta$ has no $3$-chain by Theorem~\ref{thm:Q30} and Lemma~\ref{lem:R}. On $\OG$,
$\theta$ is an antichain by Lemma~\ref{lem:og-corners}. If $C$ is a longest chain of
$P_\kappa$, then $C\setminus\theta\subseteq D(\rho)$ is a chain of the claimed size. For the
primed version, take $C$ inside the rows $\ge1$, which form copies of $P_{\kappa'}$ and
$P_{\rho'}$ (\S\ref{ssec:shifted}). Row $0$ is the largest part of both $\kappa$ and $\rho$.
\end{proof}

\begin{lemma}[A term that drops the window]\label{lem:A}
Let $\kappa\ne\varnothing$, $r=r(\kappa)$, and for $t\in\{1,2\}$ put
$T_t=\{(a,b)\in D(\kappa):a+b\ge r-t\}$ and $D(\rho_t)=D(\kappa)\setminus T_t$.
\begin{enumerate}[label=(\roman*)]
\item $\rho_t$ is a strict partition with $\rho_t\subsetneq\kappa$. If $t\le r$, then
$r(\rho_t)\le r-t$.
\item On $\LG$, let $t^*=1$ if $r$ is odd and $t^*=2$ if $r$ is even. Then $c_{\rho_{t^*}}\ne0$.
\item On $\OG$, $c_{\rho_1}=(-1)^{|T_1|-1}$.
\item $W(\rho_{t^*})\le W(\kappa)-1$ on $\LG$, and $W(\rho_1)\le W(\kappa)-1$ on $\OG$.
\end{enumerate}
\end{lemma}

\begin{proof}
(i) $T_t$ is an up-set of $P_\kappa$, so $\rho_t$ exists by Lemma~\ref{lem:ideals}. It
contains the box realising $\max(a+b)$ (Lemma~\ref{lem:chain}), so $\rho_t\ne\kappa$. If
$t\le r$ and $\rho_t\ne\varnothing$, every box of $D(\rho_t)$ has $a+b\le r-t-1$, so
$r(\rho_t)\le r-t$ by Lemma~\ref{lem:chain}. If $\rho_t=\varnothing$, then $r(\rho_t)=0\le r-t$.
(The hypothesis $t\le r$ is needed: for $\kappa=(1)$ and $t=2$ the bound would read $0\le-1$.
We only use $t=t^*\le r$.)

(ii) By Theorem~\ref{thm:Q30} and Lemma~\ref{lem:R} it suffices to show that $T_{t^*}$ has no
$3$-chain and that each of its diagonal boxes is minimal in it. The coordinate sum strictly
increases along a chain, and $T_{t^*}$ lies on at most two antidiagonals $a+b=\text{const}$, so
there is no $3$-chain. A box of $T_{t^*}$ strictly below $(c,c)$ has a smaller coordinate sum.
For $r$ odd, $T_1$ is a single antidiagonal, so no such box exists. For $r$ even,
$T_2=\{a+b\in\{r-2,r-1\}\}$, and $2c$ is even, so $2c=r-2$ lies on the inner antidiagonal;
again nothing in $T_2$ lies below $(c,c)$.

(iii) $T_1$ is the top antidiagonal: it is nonempty, and it consists of maximal elements of
$P_\kappa$. Apply Lemma~\ref{lem:og-corners}.

(iv) On $\LG$ with $r$ odd, $r(\rho_1)\le r-1$, which is even, so
$\lceil r(\rho_1)/2\rceil\le(r-1)/2=W-1$. With $r$ even, $r(\rho_2)\le r-2$, so
$\lceil r(\rho_2)/2\rceil\le r/2-1=W-1$. On $\OG$, $r(\rho_1)\le r-1=W-1$.
\end{proof}

The parity choice in (ii) matters. For $r$ odd the two-strip $T_2$ has its diagonal box (if
any) on the \emph{outer} antidiagonal, where it can lie above $(c-1,c)\in T_2$; then $V(T_2)=0$.
A machine check of the $67$ odd-$r$ two-strips at $m\le6$ against Buch--Ravikumar's formula
finds that $42$ of them have vanishing coefficient, each one exactly in this situation.

\begin{proof}[Inductive proof of Theorem~\ref{thm:window}]
We prove by strong induction on $|\kappa|$ that $s_j(\kappa)\ge0$ for all integers $j\ge1$,
with equality iff $j\le W(\kappa)$. For $\kappa=\varnothing$ this is clear. Let
$\kappa\ne\varnothing$.

\emph{Monotonicity.} For $j\ge1$, every $\rho$ in~\eqref{eq:star} has $|\rho|<|\kappa|$, so
$s_j(\rho)\ge0$ by induction, and the coefficients in~\eqref{eq:star} are nonnegative. Hence
$s_{j+1}(\kappa)\ge s_j(\kappa)$.

\emph{Anchor.} $s_1(\kappa)=0$ by Lemma~\ref{lem:anchor}. So $s_j(\kappa)\ge0$ for all $j\ge1$.

\emph{Vanishing.} Let $1\le j\le W-1$. For every $\rho$ with $c_\rho\ne0$,
Lemma~\ref{lem:chaindrop} gives $W(\rho)\ge W-1\ge j$: on $\LG$,
$\lceil r(\rho)/2\rceil\ge\lceil(r-2)/2\rceil=W-1$, and on $\OG$, $r(\rho)\ge r-1$. So
$s_j(\rho)=0$ by induction, and~\eqref{eq:star} gives $s_{j+1}(\kappa)=s_j(\kappa)$. Starting
from $s_1=0$, this reaches $s_W=0$.

\emph{Nonvanishing.} At $j=W$, \eqref{eq:star} gives
$s_{W+1}(\kappa)\ge|c_{\rho^*}|\,s_W(\rho^*)$ with $\rho^*=\rho_{t^*}$ (on $\LG$) or $\rho_1$
(on $\OG$); the other terms are $\ge0$. By Lemma~\ref{lem:A}, $c_{\rho^*}\ne0$ and
$W(\rho^*)<W$, so $s_W(\rho^*)>0$ by induction. Monotonicity gives $s_j>0$ for all $j>W$.
\end{proof}

\begin{example}
For $\kappa=(3,2)$ on $\LG(3,6)$: the coordinate sums of $D(\kappa)$ are $0,1,2,2,3$, so
$r=4$ and $W=2$. By Example~\ref{ex:pieri32}, \eqref{eq:star} at $j=2$ reads
$s_3(3,2)=s_2(3,2)+2s_2(3,1)+s_2(2,1)+s_2(3)+s_2(2)$. The first four terms vanish, since each of
those $\kappa$ has $W=2$, and $s_2(2)=1$. So $s_3(3,2)=1$, that is, $H_{(3,2)}(-3)=-1$. Here
$\rho^*=\rho_2=(2)$ and $T_2=\{(0,2),(1,1),(1,2)\}$. At $j=3$ the same recursion gives
$s_4(3,2)=1+2\cdot2+1+3+4=13$. For $\kappa=(4,3,2,1)$ we have $r=7$ and $W=4$; the support of
$[\cO_{X_\kappa}]\cdot\cO^1$ is the single term $\rho=(4,3,2)$ with $\theta=T_1=\{(3,3)\}$ and
$c_\rho=1$.
\end{example}

%======================================================================
\section{Half-integer values}\label{sec:half}
%======================================================================
% Sources: proofs/2026-10-23-odd-zero-law.md (Lemma P, Lemma W, Lemma C, Prop H, Thm A, Thm C, Thm B);
%          proofs/2026-10-24-Z-by-hyperplane-induction.md (Theorem Z);
%          proofs/2026-10-25-crownF-strict-reciprocity.md §9 (Theorem HZ).
% Registry: q488-oc-prop-H-bridge, q488-oc-thmA-half, q488-oc-thmC-M3, q488-oc-Z-by-induction,
%           q488-oc-HZ-half-integer-law, q488-oc-odd-zero-law, q488-oc-thmB-OC (all proved).

Throughout this section $X=\LG(m,2m)$, and $s_x(\kappa)=(-1)^{|\kappa|}H_\kappa(-x)$ for
$x\in\tfrac12\Z$. Since~\eqref{eq:star} is an identity of polynomials, it holds at
half-integers as well. Write $y_i=e^{-\varepsilon_i}$.

\subsection{The anchor at \texorpdfstring{$-\tfrac12$}{-1/2} and the bridge to vector bundles}
For $1\le k\le m$, put $c(k)=s_{m-k+1}s_{m-k+2}\cdots s_m$, of length $k$. Then
$D_{c(k)}=D_{m-k+1}\circ\dots\circ D_m$, so $D_m$ acts first.

\begin{lemma}[Reduced word]\label{lem:W}
$w_\kappa=c(\kappa_\ell)\cdots c(\kappa_2)\,c(\kappa_1)$, and this word is reduced. Equivalently,
$w_\kappa=w_{\kappa'}\,c(\kappa_1)$ with lengths adding, and every letter of $w_{\kappa'}$ is
$s_j$ with $j\ge m-\kappa_1+2$.
\end{lemma}

\begin{proof}
Let $p=m-\kappa_1+1$. Applying the letters of $c(\kappa_1)$ to $\omega=(1,\dots,1)$ from right to
left, $s_m$ flips the last coordinate and $s_{m-1},\dots,s_p$ carry the $-1$ to position $p$.
So $c(\kappa_1)\omega=v_{(\kappa_1)}$. Every letter of $c(\kappa_\ell)\cdots c(\kappa_2)$ is
$s_j$ with $j\ge m-\kappa_2+1\ge p+1$, so it only moves coordinates $\ge p+1$, where
$v_{(\kappa_1)}$ is $+1$. By induction on $\ell(\kappa)$, the word maps $\omega$ to
$v_\kappa$. It has length $|\kappa|=\ell(w_\kappa)$ and lies in the coset $w_\kappa W_P$, and
the minimal coset representative is the unique element of minimal length. So the word is
reduced and equals $w_\kappa$.
\end{proof}

Let $D^B_mg=(g-y_ms_mg)/(1-y_m)$ be the type-$B$ operator for the short root $\varepsilon_m$.
Put $D'_i=D_i$ for $i<m$ and $D'_m=(1+y_m)^{-1}D^B_m$, and let
$\Phi=\prod_{i=1}^m(1+y_i)^{-1}$, an analytic germ with $\Phi(0)=2^{-m}$.

\begin{lemma}[Conjugation]\label{lem:C}
For all $i$, $D_i(e^{-\omega/2}g)=e^{-\omega/2}D'_i(g)$ and $D'_i(\Phi g)=\Phi D_i(g)$.
Hence $D_w(e^{-\omega/2}g)=e^{-\omega/2}\,\Phi\,D_w(\Phi^{-1}g)$ for every $w$.
\end{lemma}

\begin{proof}
For $i<m$, both $e^{-\omega/2}$ and $\Phi$ are $s_i$-invariant, and $D_i(\varphi g)=\varphi D_ig$
whenever $s_i\varphi=\varphi$. For $i=m$: $s_m(e^{-\omega/2})=e^{-\omega/2}e^{\varepsilon_m}$ and
$e^{-\alpha_m}=y_m^2$, so
$D_m(e^{-\omega/2}g)=e^{-\omega/2}(g-y_ms_mg)/(1-y_m^2)=e^{-\omega/2}(1+y_m)^{-1}D^B_mg$.
For the second identity, write $\Phi=\Phi'(1+y_m)^{-1}$ with $\Phi'$ $s_m$-invariant, and use
$s_m(1+y_m)^{-1}=y_m(1+y_m)^{-1}$. Then
$D^B_m((1+y_m)^{-1}g)=(g-y_m^2s_mg)/((1+y_m)(1-y_m))$, and multiplying by $(1+y_m)^{-1}\Phi'$
gives $\Phi(g-y_m^2s_mg)/(1-y_m^2)=\Phi D_mg$.
\end{proof}

\begin{proposition}[Bridge]\label{prop:bridge}
For every integer $a\ge0$,
$2^mH_\kappa(-a-\tfrac12)=\bigl[D_{w_\kappa}(F_a)\bigr]_{y=1}$, where
$F_a=\prod_{i=1}^my_i^a(1+y_i)$.
\end{proposition}

\begin{proof}
By Lemma~\ref{lem:P},
\[ H_\kappa(-a-\tfrac12)=\bigl[D_{w_\kappa}(e^{-\omega/2}e^{-a\omega})\bigr](0). \]
By Lemma~\ref{lem:C} this equals $[e^{-\omega/2}\Phi\,D_{w_\kappa}(F_a)](0)$, because
$\Phi^{-1}e^{-a\omega}=F_a$, and $e^{-\omega/2}\Phi=2^{-m}$ at $0$.
\end{proof}

\begin{remark}\label{rem:wedgeU}
$F_a$ is the character of $\C_{-a\omega}\otimes\bigwedge^\bullet U_0$, where $U_0$ is the
$P$-module with weights $-\varepsilon_i$, the fibre of the tautological subbundle $U$. Pulling
back to $G/B$ and filtering by $B$-weights, (G1) turns Proposition~\ref{prop:bridge} into
$2^mH_\kappa(-a-\tfrac12)=\chi(X_\kappa,\bigwedge^\bullet U\otimes\cO(-a))$. So the half-integer
values are Euler characteristics of honest vector bundles. (For $m=1$, $X=\mathbb{P}^1$ and
$U=\cO(-1)$: $\chi(\cO\oplus\cO(-1))=1=2H(-\tfrac12)$.) None of the proofs below use this
geometric reading.
\end{remark}

\begin{theorem}\label{thm:A}
$H_\kappa(-\tfrac12)=2^{-\kappa_1}$ for every $\kappa$ (with $\kappa_1=0$ for $\kappa=\varnothing$).
\end{theorem}

\begin{proof}
By Proposition~\ref{prop:bridge} with $a=0$, $2^mH_\kappa(-\tfrac12)=[D_{w_\kappa}\prod_i(1+y_i)]_{y=1}$.
Let $p=m-\kappa_1+1$. We claim that $D_{c(\kappa_1)}\prod_{i=1}^m(1+y_i)=\prod_{i<p}(1+y_i)$.
Apply $D_m,D_{m-1},\dots,D_p$ in turn. Before $D_q$ acts, the function is
$\prod_{i\le q}(1+y_i)$ (true for $q=m$), and $\prod_{i<q}(1+y_i)$ is $s_q$-invariant. So
$D_q\prod_{i\le q}(1+y_i)=\prod_{i<q}(1+y_i)\,(1+D_qe^{-\varepsilon_q})$. Here
$\langle-\varepsilon_q,\alpha_q^\vee\rangle=-1$ for $q<m$ and for $q=m$, so
$D_qe^{-\varepsilon_q}=0$. By Lemma~\ref{lem:W}, the remaining letters $s_j$ have $j\ge p+1$; they
fix $\prod_{i<p}(1+y_i)$, and $D_j(1)=1$. Evaluating at $y=1$ gives $2^{p-1}=2^{m-\kappa_1}$.
\end{proof}

\begin{theorem}\label{thm:C}
If $\kappa_1\ge2$, then $H_\kappa(-\tfrac32)=2^{-\kappa_1}H_{\kappa'}(-1)$, and
$H_{(1)}(-\tfrac32)=-\tfrac12$. Hence $H_\kappa(-\tfrac32)=0$ iff $\ell(\kappa)\ge2$.
\end{theorem}

\begin{proof}
For $\kappa=(1)$, $H_{(1)}(j)=j+1$. Let $\kappa_1\ge2$, $p=m-\kappa_1+1$, and $f(y)=y+y^2$, so
$F_1=\prod_if(y_i)$. We apply $D_{c(\kappa_1)}$.
$D_m$ kills $y_m$ (pairing $-1$) and sends $y_m^2$ to $-1$ (pairing $-2$), so
$D_mF_1=-\prod_{i<m}f(y_i)$. Next, $D_{m-1}$ kills $y_{m-1}$ and sends $y_{m-1}^2$ to
$-y_{m-1}y_m$, so $D_{m-1}(-f(y_{m-1}))=y_{m-1}y_m$. For $p\le q<m-1$, $D_q$ applied to
$f(y_q)y_{q+1}\cdots y_m$ fixes $y_qy_{q+1}\cdots y_m$ (pairing $0$) and kills
$y_q^2y_{q+1}\cdots y_m$ (pairing $-1$). Hence
$D_{c(\kappa_1)}F_1=\prod_{i<p}f(y_i)\cdot y_p\cdot(y_{p+1}\cdots y_m)$. By
Lemma~\ref{lem:W}, $D_{w_{\kappa'}}$ acts only on the coordinates $p+1,\dots,m$, where
$y_{p+1}\cdots y_m=e^{-\omega''}$ for the rank-$(\kappa_1-1)$ block. The prefactor is inert and
equals $2^{p-1}$ at $y=1$. The Demazure expression for $H_{\kappa'}$ does not depend on the rank,
by the same inert-factor argument. Hence
$2^mH_\kappa(-\tfrac32)=2^{m-\kappa_1}H_{\kappa'}(-1)$. Finally, $H_{\kappa'}(-1)=0$ iff
$\kappa'\ne\varnothing$ (Lemma~\ref{lem:anchor}), and $H_\varnothing\equiv1$.
\end{proof}

\subsection{Vanishing by hyperplane sections}

\begin{theorem}\label{thm:Z}
% Registry: q488-oc-Z-by-induction (proved given G0, G1, BR; endorsed by Clio, review uid451).
For every strict $\kappa$ and every odd $M$ with $3\le M\le2\lceil r(\kappa')/2\rceil+1$,
$H_\kappa(-M/2)=0$.
\end{theorem}

\begin{proof}
Strong induction on $|\kappa|$. If $\ell(\kappa)\le1$, then $r(\kappa')=0$ and the range is
empty. Let $\ell(\kappa)\ge2$, so $r'=r(\kappa')\ge1$. Put $b_\kappa(x)=H_\kappa(x)-H_\kappa(x-1)$.

\emph{Every term vanishes on the shrunken window.} Let $M$ be odd with
$3\le M\le2\lceil r'/2\rceil-1$, and let $c_\rho\ne0$. By Lemma~\ref{lem:chaindrop},
$r(\rho')\ge r'-2$, so
$2\lceil r(\rho')/2\rceil+1\ge2\lceil(r'-2)/2\rceil+1=2\lceil r'/2\rceil-1\ge M$. Thus $M$ lies
in the window of $\rho$, and $H_\rho(-M/2)=0$ by induction. By~\eqref{eq:hyperplane},
$b_\kappa(-M/2)=0$.

\emph{Telescoping from the anchor.} By Theorem~\ref{thm:C}, $H_\kappa(-\tfrac32)=0$. For $M$
in the range above, $b_\kappa(-M/2)=0$ says $H_\kappa(-M/2)=H_\kappa(-(M+2)/2)$. Going up
$M=3,5,\dots,2\lceil r'/2\rceil-1$ carries the zero from $-\tfrac32$ to
$-(2\lceil r'/2\rceil+1)/2$.
\end{proof}

Theorem~\ref{thm:Z} uses only the support of the Pieri rule: it needs to know that
$c_\rho\ne0$ implies that $\theta$ has no $3$-chain. This also follows from the rim condition
of~\cite[Thm.~5.5]{BR12} together with the bound of two boxes per row and column (two boxes in
strictly south-east position, or three in a row or column, would be needed for a
$3$-chain). The coefficients do not cooperate at the edge of the window: the nonzero terms of
$b_\kappa$ at $-M^*/2$, $M^*=2\lceil r'/2\rceil+1$, need not have a common sign (for
$\kappa=(3,2)$ they are $-\tfrac18$ and $+\tfrac14$). So the nonvanishing half needs the
finer argument below.

\subsection{The odd-zero law}

Let $A_o$ and $A_e$ be the matrices indexed by strict partitions with
$A_o(\kappa,\nu)=[\theta\text{ is an antichain without diagonal boxes}]$ and
$A_e(\kappa,\nu)=[\theta\text{ is an antichain}]$, for $\nu\subseteq\kappa$ and
$\theta=D(\kappa)\setminus D(\nu)$, and $0$ otherwise.

\begin{lemma}[Layers]\label{lem:layers}
\begin{enumerate}[label=(\alph*)]
\item $(A_oA_e)(\kappa,\nu)=V(D(\kappa)\setminus D(\nu))$. Hence~\eqref{eq:star} reads
$s_{x+1}=A_oA_e\,s_x$ for every $x\in\C$.
\item For $a\ge1$, $[(A_oA_e)^{a-1}A_o](\kappa,\nu)=N_{\kappa/\nu}(a)$.
\end{enumerate}
\end{lemma}

\begin{proof}
Peel values from the top, as in Lemma~\ref{lem:peel}. The preimage of the current top value is
an antichain that is an up-set of what remains. The value is odd (no diagonal boxes allowed:
$A_o$) or even ($A_e$). In (b), the values $2a-1,2a-2,\dots,1$ alternate $o,e,\dots,o$, giving
$A_o(A_eA_o)^{a-1}=(A_oA_e)^{a-1}A_o$. Part (a) is the same with the alphabet $\{e<o\}$, and
$\theta=\varnothing$ gives $1$.
\end{proof}

\begin{lemma}[The anchor collapses to single rows]\label{lem:eps}
Let $\varepsilon=A_e\,s_{1/2}$. Then $\varepsilon(\varnothing)=1$,
$\varepsilon((L))=(-1)^{L+1}2^{-L}$ for $L\ge1$, and $\varepsilon(\kappa)=0$ if
$\ell(\kappa)\ge2$.
\end{lemma}

\begin{proof}
By Theorem~\ref{thm:A}, $s_{1/2}(\nu)=(-1)^{|\nu|}2^{-\nu_1}$, and
$\varepsilon(\kappa)=\sum_{A\subseteq\max P_\kappa}s_{1/2}(\kappa\setminus A)$. If
$\ell(\kappa)\ge2$, the last box $z$ of the bottom row is maximal and not in row $0$. Toggling
$z\in A$ changes $|\nu|$ by one and leaves $\nu_1$ alone, so the terms cancel in pairs. If
$\kappa=(L)$, the only maximal box is $(0,L-1)$, and
$\varepsilon=(-1)^L2^{-L}+(-1)^{L-1}2^{-(L-1)}=(-1)^{L+1}2^{-L}$.
\end{proof}

\begin{corollary}[Row formula]\label{cor:row}
For $a\ge1$, $s_{a+1/2}(\kappa)=N_\kappa(a)+\sum_{L=1}^{\kappa_1}(-1)^{L+1}2^{-L}N_{\kappa/(L)}(a)$.
\end{corollary}

\begin{proof}
$s_{a+1/2}=(A_oA_e)^as_{1/2}=(A_oA_e)^{a-1}A_o\,\varepsilon$; apply Lemmas~\ref{lem:layers}(b)
and~\ref{lem:eps}.
\end{proof}

For example, for $\kappa=(3)$ and $a=2$:
$s_{5/2}(3)=0+\tfrac12\cdot3-\tfrac14\cdot3+\tfrac18\cdot1=\tfrac78$, and indeed
$H_{(3)}(-\tfrac52)=-\tfrac78$.

\emph{The fibre over the lower rows.} Let $n=\kappa_1$, $M=2a-1$, and $k=\kappa_2$ (with $k=0$
if $\ell(\kappa)=1$). Rows $\ge1$ of $P_\kappa$ form a copy of $P_{\kappa'}$. Row $0$ is the
chain $(0,0)<\dots<(0,n-1)$, and its cover relations to rows $\ge1$ are $(0,\beta)<(1,\beta)$
for $1\le\beta\le k$; nothing in rows $\ge1$ lies below a box of row $0$. Fix a strictly
order-preserving $g$ on rows $\ge1$ with values in $[M]$ and even diagonal values; there are
$N_{\kappa'}(a)$ of them. Put $u_\beta=g(1,\beta)$ for $1\le\beta\le k$ and $u_\beta=M+1$ (no
bound) for $k<\beta\le n-1$. For a bound vector $u$ and $1\le L\le n$, let $G_L(M;u)$ be the
number of strictly increasing $v_L<\dots<v_{n-1}$ in $[M]$ with $v_\beta<u_\beta$ (so
$G_n=1$), and let $E_0(M;u)$ be the same count for $L=0$ with $v_0$ even. Put
\[ \Phi(u)=E_0(M;u)-\sum_{L=1}^n(-\tfrac12)^LG_L(M;u). \]
Summing Corollary~\ref{cor:row} over $g$ gives
\begin{equation}\label{eq:fibre}
s_{a+1/2}(\kappa)=\sum_g\Phi(u(g)).
\end{equation}
The bounds coming from an actual $g$ satisfy
\begin{equation}\tag{H}
u_\beta\ge\beta+1\quad\text{for }1\le\beta\le k,
\end{equation}
because $u_1=g(1,1)$ is even, hence at least $2$, and row $1$ increases strictly.

Let $\Psi_L(M;u)=\sum_{L'=L}^n(-\tfrac12)^{L'-L}G_{L'}(M;u)$. Then $\Psi_n=1$ and
$\Phi=E_0+\tfrac12\Psi_1$.

\begin{lemma}[Shift]\label{lem:shift}
Let $1\le L\le n-1$ and $M\ge1$, and suppose $u_{L'}\ge2$ for $L\le L'\le n-1$. Then
$\Psi_L(M;u)=G_L(M-1;u-1)+\tfrac12\Psi_{L+1}(M-1;u-1)$.
\end{lemma}

\begin{proof}
For $L'\ge L$, let $\alpha_{L'}$ (resp.\ $\beta_{L'}$) count the fillings of boxes
$L',\dots,n-1$ with $v_{L'}=1$ (resp.\ $v_{L'}\ge2$), and put $\beta_n=1$. Then
$G_{L'}=\alpha_{L'}+\beta_{L'}$ for $L'\le n-1$. Deleting box $L'$ (value $1$) is a bijection onto
the fillings of $L'+1,\dots,n-1$ with all values $\ge2$. Its inverse puts $1$ in box $L'$, which
is allowed because $u_{L'}\ge2$ and $M\ge1$. So $\alpha_{L'}=\beta_{L'+1}$ for $L\le L'\le n-1$.
Substituting,
\[ \Psi_L=\sum_{L'=L}^{n-1}(-\tfrac12)^{L'-L}(\beta_{L'}+\beta_{L'+1})+(-\tfrac12)^{n-L}\beta_n
=\beta_L+\tfrac12\sum_{L''>L}(-\tfrac12)^{L''-L-1}\beta_{L''}. \]
Finally, $v\mapsto v-1$ identifies fillings with values in $[2,M]$ under bounds $u$ with
fillings in $[M-1]$ under bounds $u-1$, so $\beta_{L''}(M;u)=G_{L''}(M-1;u-1)$, including
$L''=n$.
\end{proof}

\begin{lemma}[One row]\label{lem:onerow}
Assume (H).
\begin{enumerate}[label=(\alph*)]
\item If $M\ge n-1$, then
$\Phi(u)=E_0(M;u)+\sum_{i=0}^{n-2}2^{-i-1}G_{1+i}(M-1-i;u-1-i)+2^{-n}\ \ge\ 2^{-n}$.
\item If $M\le n-1$, then $\Phi(u)=(-1)^{n-1-M}2^{-n}$.
\end{enumerate}
\end{lemma}

\begin{proof}
Apply Lemma~\ref{lem:shift} repeatedly. Step $i=0,1,\dots$ rewrites $\Psi_{1+i}(M-i;u-i)$. Its
hypotheses are $M-i\ge1$ and $u_{L'}-i\ge2$ for $L'\ge1+i$. The second holds by (H) for bounded
boxes ($u_{L'}-i\ge L'+1-i\ge2$), and for unbounded boxes as long as $M-i\ge1$.
(a) If $M\ge n-1$, steps $i=0,\dots,n-2$ are allowed, ending at $\Psi_n=1$, so
$\Psi_1=\sum_{i=0}^{n-2}2^{-i}G_{1+i}(M-1-i;u-1-i)+2^{-(n-1)}$. Every term is nonnegative.
(b) If $M\le n-1$, steps $i=0,\dots,M-1$ are allowed, ending at
$\Psi_{1+M}(0;u-M)=(-\tfrac12)^{n-1-M}$, since only $G_n=1$ survives an empty alphabet. The
$G$-terms vanish: $G_{1+i}(M-1-i)$ fills $n-1-i$ boxes from $M-1-i<n-1-i$ values. So
$\Psi_1=2^{-M}(-\tfrac12)^{n-1-M}$. Also $E_0=0$, because $n$ strictly increasing values that are
all $\ge2$ need $M\ge n+1$. Hence $\Phi=\tfrac12\Psi_1$.
\end{proof}

\begin{proof}[Proof of Theorem~\ref{thm:HZ}]
% Registry: q488-oc-HZ-half-integer-law (proved given G0, G1, BR + Thm A).
The value at $-\tfrac12$ is Theorem~\ref{thm:A}. Let $a\ge1$ and $M=2a-1$. If $2a\le\kappa_1$,
then $M\le n-1$, and Lemma~\ref{lem:onerow}(b) gives
$\Phi\equiv(-1)^{n-2a}2^{-n}=(-1)^{\kappa_1}2^{-\kappa_1}$ for every $g$. So
$s_{a+1/2}(\kappa)=(-1)^{\kappa_1}2^{-\kappa_1}N_{\kappa'}(a)$ by~\eqref{eq:fibre}. If
$2a\ge\kappa_1$, then $M\ge n-1$, and Lemma~\ref{lem:onerow}(a) gives $\Phi\ge2^{-n}$ for every
$g$, so $s_{a+1/2}(\kappa)\ge2^{-\kappa_1}N_{\kappa'}(a)$. (For $2a=\kappa_1$ both apply, and
they agree because $\kappa_1$ is even.) In both cases $s_{a+1/2}(\kappa)=0$ iff
$N_{\kappa'}(a)=0$. By Lemma~\ref{lem:Npos}, this happens iff $\kappa'\ne\varnothing$ and
$a\le\lceil r(\kappa')/2\rceil$. The signs follow from $H_\kappa=(-1)^{|\kappa|}s$. Each
$\Phi(u)$ lies in $2^{-n}\Z$, so $2^{\kappa_1}s_{a+1/2}(\kappa)\in\Z$.
\end{proof}

\begin{example}
Take $\kappa=(3,2)$ and $a=2$, so $M=3\ge n-1=2$. On rows $\ge1$, the chain $(1,1)<(1,2)$ with
$(1,1)$ even and values in $[3]$ admits only $g=(2,3)$, so $u=(2,3)$. Then $E_0=0$ (no
$v_1<2$ fits above an even $v_0$) and $G_1(2;(1,2))=0$. Lemma~\ref{lem:onerow}(a) gives
$\Phi=\tfrac18$, so $s_{5/2}(3,2)=\tfrac18=2^{-\kappa_1}N_{(2)}(2)$, and
$H_{(3,2)}(-\tfrac52)=-\tfrac18$. The bound in Theorem~\ref{thm:HZ}(iii) is attained.
\end{example}

\begin{remark}[Checks]
The values $(-1)^{|\kappa|}H_\kappa(-a-\tfrac12)$, computed from type-$C$ Demazure characters
and interpolated to $x=-a-\tfrac12$, agree with $\sum_g\Phi(u(g))$ for all $\kappa$ with parts at
most $m\le4$ and $a\le5$ ($125$ cases). Lemma~\ref{lem:onerow} holds pointwise in all $477$ cases
that occur. Hypothesis (H) is necessary: without it $\Phi$ can be negative (for example $u=(1)$
gives $-\tfrac14$).
\end{remark}

\subsection{Nonvanishing for large half-integers}

\begin{proof}[Proof of Corollary~\ref{cor:OC}]
% Registry: q488-oc-thmB-OC (proved via Theorem HZ, given G0, G1, BR).
Let $M=2a+1\ge2m+1$, so $a\ge m\ge\kappa_1$ and $2a\ge\kappa_1$. The parts of $\kappa'$ are at
most $m-1$, so $\max(a'+b')\le2m-4$ on $D(\kappa')$, and $r(\kappa')\le2m-3$ by
Lemma~\ref{lem:chain}. Hence $\lceil r(\kappa')/2\rceil\le m-1<a$, and
Theorem~\ref{thm:HZ} gives $s_{a+1/2}(\kappa)\ge2^{-\kappa_1}N_{\kappa'}(a)>0$ with
$2^{\kappa_1}s_{a+1/2}(\kappa)\in\Z$. Then $2^ms_{a+1/2}(\kappa)$ is an integer, and it is at
least $2^{m-\kappa_1}\ge1$.
\end{proof}

\begin{remark}[A cohomological formula]\label{rem:serre}
Corollary~\ref{cor:OC} also follows from a positive cohomological formula. That argument uses
more geometry: normality and the Cohen--Macaulay property of Schubert varieties, Kodaira-type
vanishing via Frobenius splitting, and the formula $\omega_Y=\cO_Y(-\partial Y)\otimes L_{-\rho}$
for Schubert varieties $Y\subseteq G/B$~\cite[Ch.~3]{BK05}. Let $\pi\colon Y\to X_\kappa$ be the
$P/B$-bundle. Serre duality applied to Remark~\ref{rem:wedgeU}, followed by a $B$-stable
filtration of $\pi^*(\omega_{X_\kappa}\otimes\bigwedge^\bullet U^\vee(a))$, gives, for $a\ge m$,
\[ (-1)^{|\kappa|}2^mH_\kappa(-a-\tfrac12)=\sum_{S\subseteq[m]}h^0\bigl(Y,\cO_Y(-\partial Y)\otimes L_{\lambda_S}\bigr),\qquad\lambda_S=(a-i+[i\in S])_{i=1}^m. \]
Every $\lambda_S$ is dominant for $a\ge m$, and the term $S=[m]$ is nonzero. We record this as a
remark because the proof given above needs only (G0), (G1) and (BR).
\end{remark}

%======================================================================
\section{Examples}\label{sec:examples}
%======================================================================

\subsection{The quadric \texorpdfstring{$\LG(2,4)$}{LG(2,4)} and \texorpdfstring{$\mathbb{P}^3=\OG(2,5)$}{P3 = OG(2,5)}}
Let $\kappa=(2,1)$, so $P_\kappa$ is the chain $(0,0)<(0,1)<(1,1)$, with both ends diagonal.

\emph{Positive degrees.} For $j=1$, the order-preserving maps $g\colon P_\kappa\to\{0,1,2\}$
that are even at both ends are $(0,0,0)$, $(0,0,2)$, $(0,1,2)$, $(0,2,2)$ and $(2,2,2)$. This
gives $H(1)=5=\dim V_{\mathrm{Sp}_4}(\omega_2)$, the number of Plücker coordinates of the
quadric $\LG(2,4)\subseteq\mathbb{P}^4$.

\emph{Negative integers.} $H(j)=(j+1)(j+2)(2j+3)/6$ gives $H(-1)=H(-2)=0$, $H(-3)=-1$ and
$H(-4)=-5$, matching Example~\ref{ex:intro21}. Here $r=3$ and $W=2$. By Serre duality with
$\omega_{Q^3}=\cO(-3)$, $H(-3)=-h^0(\cO)$.

\emph{Half-integers.} $H(-\tfrac12)=\tfrac14=2^{-\kappa_1}$, and $H(-\tfrac32)=0$ because
$\ell(\kappa)=2$ (Theorem~\ref{thm:C}). Next, $H(-\tfrac52)=-\tfrac14$: here $\kappa'=(1)$ and
$r(\kappa')=1$, so the half-integer window is $a\le1$. Since $2a=4\ge\kappa_1$, the sign is
$(-1)^{|\kappa|}=-1$, and $|H(-\tfrac52)|=2^{-2}N_{(1)}(2)=\tfrac14$.

\emph{The orthogonal twin.} For $\OG(2,5)\cong\mathbb{P}^3$ with the spinor bundle
$\cO_{\mathbb{P}^3}(1)$, $\HOG(j)=\Om(P_{(2,1)},j+1)=\binom{j+3}{3}$. Its zeros are
$-1,-2,-3$, so $W=r=3$, and $\HOG(-4)=-1=-\Omb(P_{(2,1)},3)$. The window on $\OG$ is twice as
long, as Theorem~\ref{thm:window} predicts.

\subsection{All Schubert varieties of \texorpdfstring{$\LG(3,6)$}{LG(3,6)}}
Table~\ref{tab:lg3} lists $H_\kappa$ for all $\kappa\ne\varnothing$ with parts at most $3$,
together with $r=r(\kappa)$ and $r'=r(\kappa')$. Table~\ref{tab:lg3vals} lists the values at
$j=-1,\dots,-5$ and at $j=-\tfrac12,\dots,-\tfrac92$. The polynomials were computed from
Demazure characters (G1) by interpolation, independently of the theorems. The integer values
agree with the strict counts of Theorem~\ref{thm:main}(b) for $j\le4$. The zeros are exactly
those predicted by Theorems~\ref{thm:window} and~\ref{thm:HZ}: integer zeros $-1,\dots,-\lceil
r/2\rceil$, and half-integer zeros $-a-\tfrac12$ for $1\le a\le\lceil r'/2\rceil$ when
$\kappa'\ne\varnothing$.

\begin{table}[ht]
\centering
\renewcommand{\arraystretch}{1.15}
\begin{tabular}{@{}lccl@{}}
\toprule
$\kappa$ & $r$ & $r'$ & $H_\kappa(j)$\\
\midrule
$(1)$ & $1$ & $0$ & $j+1$\\
$(2)$ & $2$ & $0$ & $(j+1)^2$\\
$(3)$ & $3$ & $0$ & $(j+1)(j+2)(4j+3)/6$\\
$(2,1)$ & $3$ & $1$ & $(j+1)(j+2)(2j+3)/6$\\
$(3,1)$ & $3$ & $1$ & $(j+1)^2(j+2)(2j+3)/6$\\
$(3,2)$ & $4$ & $2$ & $(j+1)(j+2)(2j+3)(2j^2+6j+5)/30$\\
$(3,2,1)$ & $5$ & $3$ & $(j+1)(j+2)^2(j+3)(2j+3)(2j+5)/180$\\
\bottomrule
\end{tabular}
\caption{Hilbert polynomials of the Schubert varieties of $\LG(3,6)$.}\label{tab:lg3}
\end{table}

\begin{table}[ht]
\centering
\small
\renewcommand{\arraystretch}{1.2}
\setlength{\tabcolsep}{4pt}
\begin{tabular}{@{}l|rrrrr|rrrrr@{}}
\toprule
$\kappa$ & $-1$ & $-2$ & $-3$ & $-4$ & $-5$ & $-\frac12$ & $-\frac32$ & $-\frac52$ & $-\frac72$ & $-\frac92$\\
\midrule
$(1)$ & $0$ & $-1$ & $-2$ & $-3$ & $-4$ & $\frac12$ & $-\frac12$ & $-\frac32$ & $-\frac52$ & $-\frac72$\\
$(2)$ & $0$ & $1$ & $4$ & $9$ & $16$ & $\frac14$ & $\frac14$ & $\frac94$ & $\frac{25}4$ & $\frac{49}4$\\
$(3)$ & $0$ & $0$ & $-3$ & $-13$ & $-34$ & $\frac18$ & $\frac18$ & $-\frac78$ & $-\frac{55}8$ & $-\frac{175}8$\\
$(2,1)$ & $0$ & $0$ & $-1$ & $-5$ & $-14$ & $\frac14$ & $0$ & $-\frac14$ & $-\frac52$ & $-\frac{35}4$\\
$(3,1)$ & $0$ & $0$ & $2$ & $15$ & $56$ & $\frac18$ & $0$ & $\frac38$ & $\frac{25}4$ & $\frac{245}8$\\
$(3,2)$ & $0$ & $0$ & $-1$ & $-13$ & $-70$ & $\frac18$ & $0$ & $-\frac18$ & $-\frac{17}4$ & $-\frac{259}8$\\
$(3,2,1)$ & $0$ & $0$ & $0$ & $1$ & $14$ & $\frac18$ & $0$ & $0$ & $\frac18$ & $\frac{35}8$\\
\bottomrule
\end{tabular}
\caption{Values of $H_\kappa$ at negative integers and half-integers on $\LG(3,6)$.}\label{tab:lg3vals}
\end{table}

A few entries, read through the theorems:
\begin{itemize}
\item $(3)$: $s_3(3)=3$ counts the triples $2\le x<y<z\le5$ with $x$ even, namely $x=2$ and
$\{y,z\}\subseteq\{3,4,5\}$. At $-\tfrac32$ we have $2a=2\le\kappa_1=3$, so
$H=2^{-3}N_\varnothing=\tfrac18$ with sign $(-1)^{|\kappa|-\kappa_1}=+$.
\item $(3,1)$: $H(-\tfrac52)=\tfrac38>0$. Here $2a=4\ge3$, so the sign is
$(-1)^{|\kappa|}=+$, and the lower bound $2^{-3}N_{(1)}(2)=\tfrac18$ is not attained.
\item $(3,2,1)=\LG(3,6)$: $r=5$, so the integer zeros are $-1,-2,-3$. Since $\kappa'=(2,1)$ has
$r'=3$, the half-integer zeros are $-\tfrac32$ and $-\tfrac52$. Then $H(-4)=1$ and
$H(-\tfrac72)=\tfrac18=2^{-3}N_{(2,1)}(3)$.
\end{itemize}

\subsection{Computer checks}\label{ssec:checks}
Apart from the checks recorded with each section, Theorem~\ref{thm:main}(b) and
Theorem~\ref{thm:og} were compared with Hilbert polynomials computed directly from Demazure
characters at negative weights, for all $\kappa$ with parts at most $m$: $16,32,64,128$
pairs $(\kappa,j)$ with $j\le4$ for $m=2,3,4,5$, on both $\LG$ and $\OG$, with no failures.
Negative controls fail as they should (alphabet $[2j]$ instead of $[2j-1]$; no diagonal parity;
$\OG$ alphabet $[j]$). An independent census by Clio Vega, who evaluated the Demazure characters
directly at negative weights, confirmed Theorem~\ref{thm:window} for all strict $\kappa$ with
parts at most $6$, on both $\LG$ and $\OG$.

%======================================================================
\section{Open questions}\label{sec:open}
%======================================================================

\begin{question}[Other cominuscule spaces]\label{q:cominuscule}
The cominuscule $G/P$ that are not minuscule are $\LG(m,2m)$ and the odd quadrics
$Q^{2m-1}=B_m/P_1$. For the odd quadric the Bruhat poset is a chain, and exactly one cover has
bond $2$. This is the lecture-hall setting of~\cite{CCFL}. Is the Hilbert function of every
Schubert variety in a cominuscule $G/P$ the Ehrhart polynomial of its order polytope with
respect to the lattice determined by the bonds, as in Theorem~\ref{thm:main}(c)? Is there a
uniform reason, independent of the case-by-case bond computation of Lemma~\ref{lem:bonds}?
\end{question}

\begin{question}[Canonical modules]
For $j\ge1$, Kodaira-type vanishing on Schubert varieties gives
$(-1)^{|\kappa|}H_\kappa(-j)=h^0(X_\kappa,\omega_{X_\kappa}(j))$. Is there a
standard-monomial-type basis of $H^0(X_\kappa,\omega_{X_\kappa}(j))$ indexed by the strictly
order-preserving maps of Theorem~\ref{thm:main}(b)? Which $X_\kappa$ are Gorenstein, in terms of
the interior of $(\ord(P_\kappa),\Lambda)$?
\end{question}

\begin{conjecture}[A half-step operator]\label{conj:T}
Let $T$ be the unipotent square root of $A_oA_e$ (the matrix of~\eqref{eq:star}), indexed by
strict partitions ordered by inclusion. Then:
\begin{enumerate}[label=(\alph*)]
\item $s_{x+1/2}=T\,s_x$ for every $x$; equivalently, $T\,s_0=s_{1/2}$;
\item each entry $T(\kappa,\nu)$ is $0$ or $\pm2^{-k}$. It is multiplicative over the connected
components of $\theta=D(\kappa)\setminus D(\nu)$. An off-diagonal single box has weight $1$; an
off-diagonal component with at least two boxes has weight $0$; and a component containing exactly
one diagonal box is a ribbon $D$ of weight $\pm2^{-|D|}$, with sign $+$ iff the left neighbour of
its north-east end box lies in $D$.
\end{enumerate}
\end{conjecture}

Part (b) was checked on all $1716$ pairs at $m\le6$, and part (a) for $m\le4$. Neither part is
used in this paper. A natural guess, that $T$ is given by a shifted Murnaghan--Nakayama rule,
is false: the two agree only on single boxes, and their supports differ.

\begin{question}[Orthogonal half-integers]
On $\OG(m,2m+1)$ the half-integer values of $\HOG_\kappa$ have sporadic zeros, for example
$\HOG_{(3,2)}(-\tfrac52)=0$ and $\HOG_{(5,1)}(-\tfrac32)=0$. We found ten such $\kappa$ at
$m=6$, and no pattern. Is there an $\OG$ analogue of Theorem~\ref{thm:HZ}?
\end{question}

\begin{question}[Semi-toric degenerations]
Chirivì's semi-toric degeneration of $X_\kappa$~\cite{Chi00}, as summarised in~\cite[\S7]{CCFL},
has one toric component for each maximal chain $\varnothing<\dots<\kappa$, that is, for each
standard shifted tableau of shape $\kappa$. The bond vector of such a chain lies in
$\{1,2\}^{|\kappa|}$, with $1$ exactly at diagonal boxes. The associated lecture hall polytope
appears to be the simplex of Stanley's canonical triangulation of $\ord(P_\kappa)$ that
corresponds to the chain, taken in the lattice $\Lambda$. Is Chirivì's degeneration a geometric
realisation of this triangulation of $(\ord(P_\kappa),\Lambda)$? We have not compared it with
Chirivì's gluing data.
\end{question}

\begin{question}[Higher Pieri coefficients]
Proposition~\ref{prop:BRp1} is the case $p=1$ of the KLG-tableau formula~\cite[Cor.~5.6]{BR12},
whose proof is left to the reader there. A complete proof for all $p$, by the same induction
over Theorem~5.5 of~\cite{BR12}, would be useful. Do the coefficients
$c^\nu_{\lambda,p}$ admit an interpretation by strict maps into a $2p$-letter alphabet
$\{e_1<o_1<\dots<e_p<o_p\}$ with even letters on the diagonal, extending
Theorem~\ref{thm:Q30}?
\end{question}

\section*{Acknowledgements}
Clio Vega refereed the core arguments of this paper. She endorsed the hyperplane-section proof
of Theorem~\ref{thm:Z} (review uid451), at which point the Chevalley input (G0) was made
explicit at her request. She endorsed the inductive proof of Theorem~\ref{thm:window} for $\LG$
and $\OG$ (review uid457), found the missing hypothesis $t\le r$ in Lemma~\ref{lem:A}(i), and
supplied the short proof of Lemma~\ref{lem:A}(ii) from Theorem~\ref{thm:Q30}. She also replaced
our interpolation-based checks at $m=5$, which had not finished, by a direct evaluation of
Demazure characters at negative weights, extending the verification to $m\le6$. This draft was
prepared for Robin Langer's publication queue.

\bibliographystyle{amsplain}
\bibliography{refs}

\end{document}
